%!TEX root = ../thesis.tex
%*******************************************************************************
%****************************** Second Chapter *********************************
%*******************************************************************************

\chapter{Monte Carlo optimistic policy iteration}
\label{ch: mcopi}

\ifpdf
    \graphicspath{{Chapter2/Figs/Raster/}{Chapter2/Figs/PDF/}{Chapter2/Figs/}}
\else
    \graphicspath{{Chapter2/Figs/Vector/}{Chapter2/Figs/}}
\fi



\vfill

% Model-free reinforcement learning algorithms optimize policies by maximizing a value function through policy iteration, which alternates between policy evaluation and policy improvement.

% \mytodo{rewrite, idea of first paragraph is abruptly cut}

Policy iteration (PI) is a fundamental framework in reinforcement learning as it can find the optimal strategy for an MDP without knowing its transition or reward function.
Instead, PI 
learns the impact of each action on future rewards through trial and error,
stores this information in a value function
and later derives a policy with high expected returns from this value function.


Monte Carlo optimistic policy iteration (MC-O-PI) is a PI-based algorithm that iteratively refines the policy $\pol$ and the estimate $q$ by alternating between \textit{Monte Carlo policy evaluation}, which simulates complete episodes of the MDP to update $q$ towards the action values of $\pol$, 
and \textit{optimistic policy improvement}, which updates $\pol$ to obtain higher returns as anticipated by $q$.
This process enables learning the optimal policy $\best{\pol}$ and its action values $\qopt$ in unknown environments.

This chapter introduces the stochastic difference inclusion that underpins MC-O-PI, its key properties and existing convergence results.


% without relying on the transition and reward functions of the underlying MDP. cannot directly solve the Bellman equation


%%%%%%%%%%%%%%%%%%%%%%%%%%%%%%%%%%%%%%%%%
%%%%%%%%%%%%%%%%%%%%%%%%%%%%%%%%%%%%%%%%%
%%%%%%%%%%%%%%%%%%%%%%%%%%%%%%%%%%%%%%%%%
\newpage
\section{Stochastic difference inclusion}
\label{sec: diff inclusion}

We study the version of MC-O-PI where a policy improvement step follows each policy evaluation step.
These alternating steps can be represented using a single stochastic difference inclusion.
To build a clear foundation, we first discuss each step separately before combining them into a difference inclusion.


% underpinning MC-O-PI integrates the policy evaluation and improvement steps into a single update inclusion.

% Furthermore, we focus on the algorithm's asymptotic behaviour. Therefore, we assume that the target policy $\pol$


%%%%%%%%%%%%%%%%%%%%%%%%%%%%%%%%%%%%%%%%%
%%%%%%%%%%%%%%%%%%%%%%%%%%%%%%%%%%%%%%%%%
\subsection{Monte Carlo policy evaluation}

A step of Monte Carlo (MC) policy evaluation updates an estimate $q$ for the action values of a given policy $\pol$ in three steps:
\begin{enumerate}[label=\arabic*., font=\bfseries\sffamily]
    \item \textbf{\sffamily Generate episode.}
    Sample an initial state-action pair $(s_0, a_0)$, and simulate an episode following policy $\pol$:
    \begin{align*}
        (s_t, a_t, r_t)_{0 \leq t < T}
    \end{align*}
    where $a_t \sim \pol(\cdot \mid s_t)$, $r_{t} = r(s_t, a_t)$, and $s_{t+1} \sim \transit(\cdot \mid s_t, a_t)$.

    \item \textbf{\sffamily Compute returns.}
    For each pair $(s_t, a_t)$ in the episode, the observed return is related to the expected return $q_{\pi}(s_t,a_t)$ as
    \begin{align*}
        r_t + \discount r_{t+1} + \dots + \discount^{T-1-t} r_{T-1}
        =
        q_{\pi}(s_t,a_t) + v(s_t,a_t)
    \end{align*}
    for some stochastic perturbation $v(s_t, a_t)$.

    \item \textbf{\sffamily Update estimate.}
    Move $q$ towards the observed return using a stochastic approximation rule with learning rate $\alpha \in (0,1)$:
    \begin{align*}
        q_+(s_t, a_t) = q(s_t, a_t) + \lrate ( q_{\pi}(s_t,a_t) + v(s_t,a_t) - q(s_t, a_t) ) .
    \end{align*}
    
\end{enumerate}


%%%%%%%%%%%%%%%%%%%%%%%%%%%%%%%%%%%%%%%%%
%%%%%%%%%%%%%%%%%%%%%%%%%%%%%%%%%%%%%%%%%
\subsubsection{Update functions}

% The state-action pairs are thus updated asynchronously as not all pairs are updated after each episode.
% Besides, some pairs might be updated more than once with different observed returns if they were visited several times during the episode.

Not all state-action pairs are necessarily updated after the episode. There are three common methods to decide which pairs to modify: the \textit{initial}-, \textit{first}-, and \textit{every-visit} update functions.

\begin{definition}[Initial-visit updates]
\label{def: initial visit updates}
The initial-visit method updates the estimate $q$ only in the first state-action pair of an episode.
Thus given an episode $(s_t, a_t, r_t)_{0 \le t < T}$, this method only updates $q(s_0, a_0)$.
    % Thus given an episode $(s_t, a_t, r_t)_{0 \leq t < T}$,
    % this method only updates $q(s_0, a_0)$.
\end{definition}
\vspace{-0.3cm}

\begin{definition}[First-visit updates]
The first-visit method updates the estimate $q$ in each state-action pair of an episode, but only uses the observed return from the first visit to that pair.
Thus given an episode $(s_t, a_t, r_t)_{0 \le t < T}$, this method only updates $q(s_t,a_t)$ if $(s_t,a_t) \notin \{(s_{t'},a_{t'}) : t' < t\}$.
    % Thus given an episode $(s_t, a_t, r_t)_{0 \leq t < T}$,
    % this method updates $q(s_t, a_t)$ using only the observed return from the first visit to that state-action pair.
    % provided that pair does not appear in the sequence $(s_{t'}, a_{t'})_{t' < t}$.
\end{definition}
\vspace{-0.3cm}

\begin{definition}[Every-visit updates]
The every-visit method updates the estimate $q$ in each state-action pair of an episode using every observed return from that pair.
Thus given an episode $(s_t, a_t, r_t)_{0 \le t < T}$, this method updates $q(s_t,a_t)$ each time $(s_t,a_t)$ was visited during the episode.
    % Thus given an episode $(s_t, a_t, r_t)_{0 \leq t < T}$,
    % this method updates $q(s_t, a_t)$ using every observed return from that state-action pair.
    % this method updates $q(s_t, a_t)$, each time that pair is encountered.
\end{definition}
\vspace{-1.2\parskip}


% In other words, given an episode $(s_t, a_t, r_t)_{0 \leq t < T}$, the update function determines in which state-action pairs the estimate $q$ is modified.
% The initial-visit method updates only the first state-action pair $(s_0, a_0)$.
% The first-visit method updates each state-action pair $(s_t, a_t)$ provided it has not appeared in the sequence $(s_{t'}, a_{t'})_{t' < t}$.
% The every-visit method updates all state-action pairs each time they are encountered.
We only consider initial- and first-visit updates
because every-visit updates can produce biased estimates of the action values \cite[Theorem~7]{Singh1996ReinforcementTraces}.

Under initial- and first-visit updates, each state-action pair of the estimate $q$ is updated at most once during a step of MC policy evaluation.
The update is then equivalent to
\begin{align}
\label{eq: update q with chi}
    q_+ 
    =
    q + \lrate u ( q_\pol + v - q )
    % \begin{cases}
    %     \ q(s,a) + \lrate (\Tilde{q}_\pol(s,a) - q(s,a)) &\text{if} \ u(s,a)=1
    %     \\
    %     \ q(s,a) &\text{if} \ u(s,a)=0
    % \end{cases}
\end{align}
where $u$ is an indicator function satisfying $u(s,a)=1$ if $(s, a)$ is to be updated, and $0$ otherwise.
This function depends on the initial state-action pair of the episode, the policy used for simulating it and the chosen update function.


% The major difference between initial- and first-visit updates, is the probability that a state-action pair is updated.
% For initial-visit MC-O-PI, a state-action pair can only be updated if the episode starts in that particular state-action pair.
% For first-visit MC-O-PI on the other hand, a state-action pair can also be updated if the simulation passes through it later on. 
% Therefore, for on-policy actions, we have that
% \begin{align}
%     \saprior(s, \pol(s)) = \update^{IV}(s, \pol(s)) \leq \update^{FV}(s, \pol(s)) ,
%     % \leq \update^{EV}(s, \pol(s)) ,
% \end{align}
% whereas for off-policy actions $a \notin \pol(s)$ we have that
% \begin{align}
%     \saprior(s, a) = \update^{IV}(s, a) = \update^{FV}(s, a) .
%     % = \update^{EV}(s, a).
% \end{align}
% Thus, depending on the update function, the update probabilities $\update$ might depend on the policy $\pol$.



% \textcolor{red}{How implement synchronous updates for IV and FV with biased sampling of initial (s,a)? need batch of episodes?}

% \textcolor{red}{instead of synchronous and asynchronous talk about MCES and expected MCES?}


%%%%%%%%%%%%%%%%%%%%%%%%%%%%%%%%%%%%%%%%%
%%%%%%%%%%%%%%%%%%%%%%%%%%%%%%%%%%%%%%%%%
\subsubsection{Initial sampling distribution}

To ensure that all state-action pairs are updated sufficiently often,
we commonly impose \textit{exploring starts} on the sampling distribution of the initial state-action pair of each episode.

\begin{definition}[Exploring starts]
\label{def: exploring starts}
    Every state-action pair has a non-zero probability of being selected as the start of an episode, meaning for every state $s \in \sset$ and action $a \in \aset(s)$, the distribution $\start$ of the initial state-action pair satisfies
    \begin{align*}
        \start(s,a) > 0.
    \end{align*}
\end{definition}
\vspace{-0.4cm}

We can effectively nullify this condition using quickly decaying sampling distributions.
For example, let $k$ denote the episode number. If we set the starting probability in a specific state-action pair $(s,a)$ to $\sigma_k(s,a) = 1/(k+2)^2$, there is a $50\%$ chance that no episode ever begins there. Under initial-visit updates, that pair would thus never be updated. To prevent such issues, we introduce a stronger condition.

\begin{definition}[Strict exploring starts]
\label{def: strict exploring starts}
    There exists a fixed $\sigma_{\min} \in \mathbb R_{> 0}$ such that
    for every state $s \in \sset$ and action $a \in \aset(s)$, the distribution $\start$ of the initial state-action pair satisfies
    \begin{align*}
        \start(s,a) > \sigma_{\min}.
    \end{align*}
\end{definition}
% \vspace{-1.8\parskip}

%%%%%%%%%%%%%%%%%%%%%%%%%%%%%%%%%%%%%%%%%
%%%%%%%%%%%%%%%%%%%%%%%%%%%%%%%%%%%%%%%%%
\subsubsection{Update distribution}

Given the sampling distribution $\sigma$ of the initial state-action pair, the policy $\pi$ used for simulating the episode and the chosen update function $f$,
we can compute the probability of updating each state-action pair after the episode.
We gather these probabilities in the update distribution $\mu : \sset \times \aset \to [0,1]$ defined as
\[
\mu := \mathbb E [ u \mid \sigma, \pi, f] .
\]
Using this distribution, we can combine the stochastic terms $u$ and $v$ in equation \eqref{eq: update q with chi} into a single term $w$ defined as
\begin{align}
\label{eq: definition combined noise}
    w
    \coloneqq 
    u \, v + (u - \mu) ( q_{\pol} - q ) .
\end{align}
This term incorporates the perturbations caused by the episode simulation and the asynchronous updates of state-action pairs.
It simplifies equation \eqref{eq: update q with chi} to
\begin{align}
\label{eq: mcopi not yet sri}
    q_+
    =
    q + \lrate ( \update ( q_{\pol} - q ) + w ) .
\end{align}

\vspace{-0.4cm}

\begin{remark}[Mean-field ODE]
\label{rmk: approximate ODE}
As we will show later, $w$ is a zero-mean stochastic perturbation.
So MC policy evaluation is essentially a stochastic approximation of the differential equation
\[
\dot q = \update ( q_{\pol} - q ) .
\]
The update distribution $\mu$ 
% resulting form the chosen sampling distribution and update function,
is thus a crucial design parameter of the algorithm
as it controls the behaviour of this underlying ODE.
\end{remark}

\vspace{-0.3cm}

\begin{remark}[Exploration-exploitation trade-off]
\label{rmk: exploration-exploitation tradeoff}
The update distribution $\mu$ controls the exploration-exploitation trade-off
as it determines how frequently each action is updated.
Exploration favours actions whose current estimated value is most uncertain, for instance, by biasing $\mu$ towards actions that received few updates.
Exploitation, on the other hand, prioritises updating actions with the highest estimated value. It thus biases $\mu$ towards actions that are currently greedy with respect to $q$.

Common strategies implement different choices of $\mu$. 
For instance, the epsilon-greedy approach starts with a uniform distribution that progressively concentrates on greedy actions as $\varepsilon$ decays.
The UCB method 
initially favours
% assigns higher probabilities to 
actions with fewer past updates and gradually shifts to greedy actions as uncertainty decreases.
\end{remark}


%%%%%%%%%%%%%%%%%%%%%%%%%%%%%%%%%%%%%%%%%
%%%%%%%%%%%%%%%%%%%%%%%%%%%%%%%%%%%%%%%%%
\subsection{Optimistic policy improvement}

Given an estimate $q$ of the action values $q_\pol$, optimistic policy improvement chooses a new policy $\pol_+$ so that, for every state $s \in \sset$,
\begin{align}
\label{eq: optim policy improvement idea}
    q(s, \pol_+) \geq q(s, \pol) .
\end{align}
The improvement is \textit{optimistic} because the policy is updated before $q$ has converged to $q_{\pol}$. Therefore, the \hyperref[thm: policy improvement theorem]{Policy Improvement Theorem} does not guarantee that $\pol_+$ is as good as $\pol$.

There are several ways to compute a policy $\pol_+$ that satisfies condition \eqref{eq: optim policy improvement idea}.
Yet, the GLIE conditions (\autoref{def: glie}) require that the policy $\pol$ is asymptotically greedy with respect to the estimate $q$, so that solutions stop moving once they reach the optimal policy $\best{\pol}$ and its action values $\qopt$.
Since we study the asymptotic behaviour of MC-O-PI, we assume policies are greedy at every step.

Concretely, we define the set-valued function $\grpol : \qset \rightrightarrows \polset$ that maps any action-value function $q$ to the set of deterministic policies that are greedy with respect to $q$:
\begin{align}
\label{eq: greedy pol}
    \grpol(q) \coloneqq \big\{ \pol \in \polset : q(s, \pol) = \max_{a \in \aset(s)} q(s,a), \ \forall s \in \sset \big\} ,
\end{align}
and choose
\begin{align}
\label{eq: greedy policy update}
    \pol_+ \in \grpol(q)
\end{align}
using random tie-breaking that is independent across states. Specifically, we assume that there exists a fixed $\beta>0$ such that, for every state $s \in \sset$ and every greedy action $a \in \arg\max_{\aset(s)} q_k(s,\cdot)$,
\begin{equation}
\label{eq:greedy-pol-tie-breaking}
    \mathbb P(\pi_+(s) = a ) \ge \beta.
\end{equation}
This assumption is mainly technical as ties rarely occur in practice, except perhaps at initialization.


%%%%%%%%%%%%%%%%%%%%%%%%%%%%%%%%%%%%%%%%%
%%%%%%%%%%%%%%%%%%%%%%%%%%%%%%%%%%%%%%%%%
\subsection{Difference inclusion combining evaluation and improvement steps}

MC-O-PI thus learns the optimal policy of an MDP by tracking two variables: an estimate $q$ and a policy $\pol$. The algorithm alternates between 
Monte Carlo policy evaluation that updates $q$ towards the action values of $\pol$ using equation \eqref{eq: mcopi not yet sri}, and optimistic policy improvement that updates $\pol$ to be greedy with respect to $q$ using equation \eqref{eq: greedy policy update}.
We can combine both equations into the following difference inclusion.

\begin{definition}[MC-O-PI difference inclusion]
\label{def: mcopi}
Given a sequence of learning rates $(\lrate_k)_{k \geq 0}$, a sequence of start distributions $(\start_k)_{k \geq 0}$ and a sequence of update functions $(f_k)_{k \geq 0}$,
MC-O-PI is equivalent to the stochastic difference inclusion
\begin{align}
\nonumber
    &q_{k+1}
    \in
    \Big\{ 
    q_k + \lrate_k \bigl( \update_k ( q_{\pol_k} - q_k ) + w_k \bigr) 
    :
    \\
\label{eq: complete difference inclusion}
    &\hspace{1cm}
    \pol_k \in \grpol(q_k),
    \update_k = \update(\start_k, \pol_k, f_k),
    w_k \sim w(\start_k, \pol_k, f_k)
    \Big\}
\end{align}
where $\update_k$ is the update distribution obtained by sampling the initial state-action pair of the $k^\nth$ episode according to $\start_k$, simulating the episode with policy $\pol_k$ and deciding which state-action pairs to update using the update function $f_k$.
The stochastic perturbation $w_k$ is a realisation under the same conditions.
\end{definition}

\vspace{-0.6cm}

System \eqref{eq: complete difference inclusion} is a switched linear system whose active dynamics depend on which policy is greedy with respect to the current estimate $q_k$.
Thus, the dynamics remain the same as long as the greedy policy does not change.
We can delimit the subspaces of $\qset$ with constant dynamics using the set-valued function $\graval : \polset \rightrightarrows \qset$ defined as
\begin{align}
\label{eq: greedy action values}
    \graval(\pol) \coloneqq \big\{ q \in \qset : q(s, \pol) = \max_{a \in \aset(s)} q(s,a), \ \forall s \in \sset \big\} .
\end{align}
This function is the inverse of $P$ defined in \eqref{eq: greedy pol}, as it maps any deterministic policy $\pol$ to the set of action-value functions for which $\pol$ is greedy.
Since $\graval(\pol)$ forms a convex subset of $\qset$, we refer to $\graval(\pol)$ as the \textit{region} of policy $\pol$ in $\qset$.


For instance, consider the MDP in \autoref{fig: description region smallest MDP} with $\sset = \{s\}$ and $\aset(s) = \{a,a'\}$.
Any point $q \in \qset$ can be represented by a point in $\real^2$ with coordinates $( q(s,a), q(s,a') )$.
As shown in \autoref{fig: show regions}, the line $q(s,a) = q(s,a')$ bisects $\qset$.
The region $\graval(\pol)$ is the half-space where $q(s,a) \ge q(s,a')$, while $\graval(\pol')$ is the half-space where $q(s,a') \ge q(s,a)$.

\label{page: description of switched system}

Suppose the update distributions $(\mu_k)$ are uniform over all state-action pairs, and hence equivalent to a scalar in \eqref{eq: complete difference inclusion}.
Following the ODE-interpretation in \autoref{rmk: approximate ODE},
MC-O-PI solutions then approximate $\dot q = q_{\pi} - q$ over $\graval(\pi)$ and $\dot q = q_{\pi'} - q$ over $\graval(\pi')$, or in short
\begin{align}
\label{eq: switched ode for interpretation}
\dot q \in \{ q_\pi - q : \pi \in \grpol(q) \}.
% \begin{cases}
%     &q_{\pi} - q \qquad q \in \graval(\pi)
%     \\
%     &q_{\pi'} - q \qquad q \in \graval(\pi')
% \end{cases}
\end{align}
As the learning rate decreases, MC-O-PI follows the streamlines in \autoref{fig: explain regions with streamlines} more closely (arbitrary position of $q_\pi$ and $q_{\pi'}$).
On the boundary, either systems are possible since the boundary belongs to both regions. However, the tie-breaking must satisfy \eqref{eq:greedy-pol-tie-breaking}.


% Note that the function $\graval$ does not partition the space $\qset$ because regions of different policies can overlap.
% For instance, if $q(s, \pol) = q(s, \pol') = \max_{a} q(s,a)$ for all states $s \in \sset$,
% then $q$ lies on the boundary between the regions $\graval(\pol)$ and $\graval(\pol')$.


\begin{figure}[hb!]
    \centering
    \hrule
    \begin{subfigure}[c]{\textwidth}
        \centering
        \vspace{0.6cm}
        \begin{tikzpicture}[
            scale=1, auto, node distance=2cm, 
            every edge/.style={draw, <-, font=\small},
            state/.style={circle, fill=mylightgrey, draw=black, text=black}
        ]
            % Nodes
            \node[state] (s1) at (0,0) {$s$};
            % Edges
            \path[->]
                (s1) edge[out=160, in=200, loop, looseness=8, draw=black] node[pos=0.5, left, text=black] {\textcolor{mygrey}{$\pi(s) :=$} $a$} (s1)
                (s1) edge[out=20, in=340, loop, looseness=8, draw=black] node[pos=0.5, right, text=black] {$a'$ \textcolor{mygrey}{$ =: \pi'(s)$}} (s1) ;
        \end{tikzpicture}
        \vspace{0.3cm}
        \caption{MDP}
        \label{fig: description region smallest MDP}
    \end{subfigure}%
    % 
    \\[3ex]
    \pgfmathsetmacro{\boxsize}{2.4}
    \begin{subfigure}[t]{0.48\textwidth}
        \centering
        \begin{tikzpicture}[scale=1.5]

                % \filldraw[fill=mylightgrey, draw=none] (0,0) rectangle (\boxsize,\boxsize);
                
                \draw[->] (0,0) -- (\boxsize,0) node[below] {$q(s,a)$};
                \draw[->] (0,0) -- (0,\boxsize) node[above] {$q(s,a')$};

                \draw[thick, CambridgeCoreBlue] (0,0) -- (\boxsize,\boxsize);

                \node at (\boxsize/2,\boxsize/2) [font=\small, color=CambridgeCoreBlue, rotate=45, anchor=north] {$q(s,a) = q(s,a')$};
                
                \node at (1.8, 0.6) [font=\small, color=black] {$\region(\pi)$};
                \node at (0.8,1.6) [font=\small, color=black] {$\region(\pi')$};
        
            \end{tikzpicture}
        \caption{Regions of policies $\pi$ and $\pi'$ \\as defined in \eqref{eq: greedy action values}}
        \label{fig: show regions}
    \end{subfigure}
    \hfill
    \begin{subfigure}[t]{0.48\textwidth}
        \centering
        \begin{tikzpicture}[scale=1.5]

            \draw[->] (0,0) -- (\boxsize,0) node[below] {$q(s,a)$};
            \draw[->] (0,0) -- (0,\boxsize) node[above] {$q(s,a')$};

            \draw[black, dotted] (0,0) -- (\boxsize,\boxsize);

            \begin{axis}[
                hide axis,
                xmin=0, xmax=\boxsize,
                ymin=0, ymax=\boxsize,
                width=\boxsize cm,
                height=\boxsize cm,
                at={(0,0)},
                scale only axis,
                clip=true,
                anchor=south west,
            ]
                \addplot[mygrey, no marks] 
                table [ x expr=\thisrowno{0}, 
                        y expr=\thisrowno{1}, 
                        col sep=comma, 
                        unbounded coords=jump] 
                {figures/algo/stream-regions-explain.csv};
            \end{axis}
            
            \filldraw[fill=black, draw=black] (2,1.4) circle (1pt) node[right, text=black, font=\small] {$q_\pol$};
            \filldraw[fill=black, draw=black] (1.2,0.4) circle (1pt) node[left, text=black, font=\small] {$q_{\pol'}$};

        \end{tikzpicture}
        \caption{Mean-field dynamics \eqref{eq: switched ode for interpretation} \\with uniform update distributions}
        \label{fig: explain regions with streamlines}
    \end{subfigure}%
    \caption{MC-O-PI is a switched linear system whose dynamics are constant over the region of each policy.}
    \label{fig: description of region of policy}
\end{figure}

\begin{remark}[Every boundary has a best and a worst policy]
\label{rmk: boundary and optimal policy}

    Consider an MDP with states $\sset$ and actions $\aset(s)$ for each $s \in \sset$. These sets define the function space $\qset \coloneqq \{ \sset \times \aset \to \real \}$.
    
    For every $q \in \qset$, we can construct a smaller MDP by removing in each state all actions from $\aset(s)$ that do not maximise $q(s,\cdot)$.
    This way, all policies of the smaller MDP are greedy with respect to $q$,
    and all these greedy policies are legal on the smaller MDP.
    The set $\grpol(q)$ thus contains all deterministic policies of this smaller MDP.
    
    The smaller MDP still satisfies the Bellman optimality principle. Therefore, it also has at least one optimal and one worst policy.
    As a result, for every $q \in \qset$ the set $\grpol(q)$ contains at least one policy that is better than or equal to all others, and another that is worse than or equal to all others.
    In particular, when $q$ lies on the boundary between several regions $\region(\pol)$, $\region(\pol')$, etc., there exist at least one best and one worst policy among those sharing this boundary.
\end{remark}


% \begin{remark}
% \label{rmk: only qopt is inside its own region}
%     Since the Bellman operator has a unique fixed point, only the optimal action-values are inside their own region, namely
%     \begin{align}
%         q_\pol \in \region(\pol) \iff \B q_\pol = \B_\pol q_\pol = q_\pol = q_{\best{\pol}}.
%     \end{align}
%     \textcolor{red}{cannot use bellman operator yet}
% \end{remark}


%%%%%%%%%%%%%%%%%%%%%%%%%%%%%%%%%%%%%%%%%
%%%%%%%%%%%%%%%%%%%%%%%%%%%%%%%%%%%%%%%%%
\subsubsection{Stochastic perturbations form a martingale difference sequence}

Recall from \eqref{eq: definition combined noise} that the perturbation in \eqref{eq: complete difference inclusion} is
\[
    w_k
    \coloneqq 
    u_k v_k + (u_k - \mu_k) ( q_{\pol_k} - q_k ) ,
\]
where $u_k$ indicates which state-action pairs are updated after the $k^\nth$ episode, $v_k$ is the error between the observed and the expected returns under policy $\pi_k$, and $\mu_k$ is the expected value of $u_k$ given the sampling distribution $\sigma_k$, the policy $\pi_k$ and the update function $f_k$.

Under initial- or first-visit updates, episode returns are unbiased estimates of $q_{\pi_k}$ and have uniformly bounded variance \cite[Theorems~6~\&~8]{Singh1996ReinforcementTraces}.
So there exists a constant $c_v \in [0,\infty)$ such that for every time $k \ge 0$, every state $s \in \mathcal S$ and every action $a \in \mathcal A(s)$,
\begin{align}
\label{eq: v_mds_first}
% &\mathbb E[v_k(s,a) \mid \mathcal F_k, \pi_k, u_k(s,a)=1] = 0 ,
&\mathbb E\!\left[v_k(s,a) \mid \mathcal F_k \right] = 0 ,
\\
\label{eq: v_mds_second}
% &\mathbb E[v_k^2(s,a) \mid \mathcal F_k, \pi_k, u_k(s,a)=1] \le c_v .
&\mathbb E\!\left[v_k^2(s,a) \mid \mathcal F_k \right] \le c_v ,
\end{align}
where $\mathcal F_k$ represents the available information after $q_k$ is computed,
but before selecting the next greedy policy, sampling the initial state-action pair, and generating the episode.
The sequence $(v_k)_{k \ge 0}$ is thus a square integrable martingale difference sequence adapted to the filtration $(\mathcal F_{k})_{k \ge 0}$.
The stochastic perturbations $(w_k)_{k \ge 0}$ inherit this property.

\begin{lemma}
[Martingale difference perturbations]
\label{thm: properties of noise}

Under initial- or first-visit updates, 
there exists a $c_w \in [0,\infty)$
such that for every time $k \ge 0$, every state $s \in \sset$ and every action $a \in \aset(s)$,
\vspace{-0.3cm}
\begin{align}
\label{eq: muw_mds}
&\mathbb E [w_k(s,a) \mid \mathcal F_k ] = 0,
\\
\label{eq: muw_second_moment}
&\mathbb E [w_k^2(s,a) \mid \mathcal F_k ] \le c_w(1+\|q_k\|^2) .
\end{align}
\end{lemma}
\vspace{-0.9cm}
\begin{proof}
    Fix a time $k\ge 0$, a state $s \in \mathcal S$ and an action $a \in \mathcal A(s)$.
    Let $\mathcal F_k'$ represent the available information right after selecting the policy but before sampling the initial state-action pair, so $\mathcal F_k \subseteq \mathcal F_k'$.
    By definition,
    \[
    w_k(s,a) = u_k(s,a) v_k(s,a) + \big( u_k(s,a)-\mu_k(s,a) \big) \big( q_{\pi_k}(s,a) - q_k(s,a) \big).
    \]
    % \textcolor{red}{wrong, $u_k$ is fixed given $F'_k$ so no expectation left, likely need $F'_k$ just right after selecting the policy but before selecting initial state-action pair}
    Since episode returns are unbiased under initial- or first-visit updates,
    and $u_k(s,a) v_k(s,a) = 0$ on $\{u_k(s,a)=0\}$,
    it follows that
    \begin{align*}
    &\mathbb E[ u_k(s,a) v_k(s,a) \mid \mathcal F_k']
    \\
    &\qquad=
    \mathbb E \big[ 
        \mathbf 1_{\{u_k(s,a)=1\}} \mathbb E[ v_{k}(s,a) \mid \mathcal F_k', u_k(s,a)=1] 
        % +
        % \mathbf 1_{\{u_k(s,a)=0\}} \mathbb E[ v_{k}(s,a) \mid \mathcal F_k', u_k(s,a)=o] 
        \mid \mathcal F_k' \big]
    = 0 .
    \end{align*}
    Also, $\mu_k(s,a) = \E[u_k(s,a) \mid \mathcal F_k']$ and,
    given that $k$, $s$ and $a$ are fixed,
    $\mu_k(s,a)$ and $(q_{\pi_k}(s,a)-q_k(s,a))$ are $\mathcal F_k'$-measurable. It follows that
    \begin{align*}
    &\mathbb E \!\left[
    \big( u_k(s,a)-\mu_k(s,a) \big) \big( q_{\pi_k}(s,a)-q_k(s,a) \big)
    \mid \mathcal F_k'
    \right]
    \\
    &\qquad =
    \big( \mathbb E [
     u_k(s,a)
    \mid \mathcal F_k'
    ]
    - \mu_k(s,a) 
     \big) 
    \big( q_{\pi_k}(s,a)-q_k(s,a) \big) 
    =
    0 .
    \end{align*}
    Thus $\mathbb E[w_k(s,a) \mid \mathcal F_k'] = 0$.
    Applying the tower property then yields
    \[
    \mathbb E [ w_{k}(s,a) \mid \mathcal F_k ]
    = \mathbb E \!\left[ \mathbb E [ w_{k}(s,a) \mid \mathcal F_k'] \mid \mathcal F_k \right]
    = 0 ,
    \]
    which proves \eqref{eq: muw_mds}.
    
    Further, since $u_k^2(s,a) \le 1$ and $(u_k(s,a) - \mu_k(s,a))^2 \le 1$,
    using $(x+y)^2 \le 2 x^2 + 2 y^2$ yields
    \[
    w_k^2(s,a) \le 2 v_k^2(s,a) + 2 (q_{\pi_k}(s,a)-q_k(s,a) )^2 .
    \]
    Similarly,
    \[
    ( q _{\pi_k}(s,a)-q_k(s,a) )^2
    \le
    2 q_{\pi_k}^2(s,a) + 2 q_k^2(s,a) .
    \]
    Since all action-values are bounded,
    there exists a constant $c_\polset < \infty$ such that $\|q_\pi\|^2 \leq c_\polset$ for every policy $\pi \in \polset$. As a result,
    \begin{align*}
    \mathbb E [ w_{k}^2(s,a) \mid \mathcal F_k ]
    &\le
    2 c_v + 2 \, \mathbb E \!\left[ ( q_{\pi_k}(s,a)-q_k(s,a) )^2 \mid \mathcal F_k \right]
    \\
    &\le
    2 c_v + 4 c_\polset + 4 \|q_k\|^2 .
    \end{align*}
    So \eqref{eq: muw_second_moment} holds for example with $c_w := \max \{ 2 c_v + 4 c_\polset, 4\}$.
\end{proof}

Since the perturbations $(w_k)_{k \ge 0}$ form a martingale difference sequence,
% adapted to the filtration $(\mathcal F_{k})$,
we can bound their cumulative impact on any solution of \eqref{eq: complete difference inclusion} using appropriate conditions on the learning rates.


%%%%%%%%%%%%%%%%%%%%%%%%%%%%%%%%%%%%%%%%%
%%%%%%%%%%%%%%%%%%%%%%%%%%%%%%%%%%%%%%%%%
\subsubsection{Modified Robbins-Monro conditions for effective learning rates}
\label{sec: modified robbins monro conditions}

% \textcolor{red}{wlog assume everywhere in the thesis that learning rates are in $(0,1)$}

Unless stated otherwise, we assume that $\lrate_k \in (0,1)$. In some examples, however, we use \textit{sample averaging} learning rates $\lrate_k(s,a) = 1/n_k(s,a)$, where $n_k(s, a)$ denotes the number of times the estimate $q$ has been updated in the state-action pair $(s, a)$ during the first $k$ time steps. 
These learning rates are theoretically appealing because they ensure that $q_k(s,a)$ is the exact average of all returns observed up to step $k$.
However, they are rarely used in practice due to their high memory requirements.
Regardless, the learning rates of stochastic approximation algorithms such as \eqref{eq: complete difference inclusion} must typically satisfy the Robbins-Monro conditions \citep{Robbins1951AMethod}.
These conditions ensure that the learning rates are large enough to avoid premature convergence due to Zeno-like behaviour,
while gradually decreasing to zero so that the accumulated stochastic perturbations
% satisfy $\lim_{K \to \infty} \sum_{k \ge K} \alpha_k w_k = 0$, and hence
cannot prevent convergence.

% \vspace{0.2cm}

\begin{definition}[Robbins–Monro conditions]
\label{asp: stochastic approximation step size}
The sequence of learning rates $(\lrate_k)_{k\geq0}$ satisfies
\begin{align}
    \label{eq: robbins-monro first}
    \sum_{k = 0}^{\infty} \lrate_\kpz &= \infty , 
    \\
    \label{eq: robbins-monro second}
    \sum_{k = 0}^{\infty} \lrate_\kpz^2 &< \infty .
\end{align}
\end{definition}

\vspace{-0.5cm}

For MC-O-PI, however, the update distributions $(\update_k)$ modulate the learning rates $(\lrate_k)$, leading to different learning rates for different state-action pairs.
Thus, Robbins-Monro-like conditions are necessary for the products $(\lrate_k \update_k)$ rather than the original learning rates alone.
We refer to these products as \textit{effective learning rates}.

To see why the Robbins-Monro conditions must be adjusted, fix a policy $\pi$ and consider the system
\begin{align}
\label{eq: single system}
    q_{k+1}
    &= q_k + \lrate_k \update_k (q_\pi - q_k) + \lrate_k w_k .
\end{align}
Unrolling over $n$ steps yields
\begin{multline*}
    q_{k+n}
    = \left( \prod_{i=k}^{k+n-1} (\I -\lrate_i \update_i) \right) q_k 
    + \left( \I - \prod_{i=k}^{k+n-1} (\I -\lrate_i \update_i) \right) q_\pi
    \\
    + \sum_{j=k}^{k+n-1} \prod_{i=j+1}^{k+n-1} (\I -\lrate_i \update_i) \lrate_j w_j .
\end{multline*}
For $q_{k+n}$ to converge to $q_\pi$ as $n$ grows to $\infty$, three conditions must be satisfied for every state $s \in \sset$ and action $a \in \aset(s)$.
First,
\begin{align}
\label{eq: condition infinite visits}
    \sum_{k=0}^{\infty} \update_k(s,a) = \infty ,
    % \qquad \forall \, s \in \sset, \forall \, a \in \aset(s) .
\end{align}
so that the state-action pair is updated infinitely often.
% Therefore, for every state-action pair $(s,a)$, the update distributions $(\update_k)$ must satisfy
Second, 
\begin{align}
    \label{eq: condition infinite effective learning rate}
    \sum_{k=0}^{\infty} \lrate_k \update_k(s,a) = \infty ,
    % \qquad \forall \, s \in \sset, \forall \, a \in \aset(s) .
\end{align}
so that the product $\prod_{i=k}^{k+n-1} (1 - \lrate_i \update_i(s,a))$ decays to zero for increasing $n$.
Finally,
\begin{align}
    \label{eq: summable learning rates outside def}
    \sum_{k=0}^\infty \lrate_k^2 < \infty ,
\end{align}
to ensure the accumulated stochastic perturbations vanish.
% Note that \eqref{eq: condition infinite visits} does not automatically imply \eqref{eq: condition infinite effective learning rate}.
% For instance, under \eqref{eq: condition infinite visits} and the Robbins-Monro conditions on the learning rates $(\lrate_k)$, we could choose $\lrate_k = \update_k(s,a) = 1/k$. System \eqref{eq: single system} would then converge prematurely due to Zeno-like behaviour because the products $\prod_{i=k}^{k+n-1} (1-\lrate_i \update_i)$ would not decay to 0.
% If we imposed $\sum_k \lrate_k^2 \update_k^2 < \infty$ instead, we could choose $\lrate_k = \update_k(s,a) = 1/\sqrt{k}$ for each state-action pair, still satisfy conditions \eqref{eq: condition infinite visits} and \eqref{eq: condition infinite effective learning rate}.
%  However, the noise term in \eqref{eq: single system unrolled} would become
% \begin{align}
%     \sum_{j=k}^{k+n-1} \prod_{i=j+1}^{k+n-1} (1-\frac{1}{i}) \frac{1}{\sqrt{j}} \xi_j
%     &= \sum_{j=k}^{k+n-1} \frac{j}{k+n+1} \frac{1}{\sqrt{j}} \xi_j
%     \\
%     &= \frac{1}{k+n+1} \sum_{j=k}^{k+n-1} \sqrt{j} \xi_j
%     \\
%     &\eqqcolon N_{k,n}
% \end{align}
% We see that
% \begin{align}
%     \E[N_{k,n}^2] = \E \Bigg[ \Big( \frac{1}{k+n+1} \sum_{j=k}^{k+n-1} \sqrt{j} \xi_j \Big)^2 \Bigg]
% \end{align}
% Since the cross terms vanish in expectation, 
% and as a result, $q_k$ would not converge to $q_\pi$
% 
% since $\sum_k \lrate_k^2 = \infty$, the noise would not be almost surely bounded. In other words, the noise alone would be able to pull solution arbitrarily far from $q_\pi$ infinitely often, thereby preventing convergence.
% 

% Thus, the minimal requirements on the learning rates and the update distributions are
% % for every state $s \in \sset$ and action $a \in \aset(s)$
% \begin{align}
%     \label{eq: xtn robbins-monro first}
%     &\sum_{k =0}^\infty \lrate_k \update_k(s,a) = \infty 
%     \qquad \forall \, s \in \sset, \forall \, a \in \aset(s) ,
%     \\
%     \label{eq: xtn robbins-monro second}
%     &\sum_{k =0}^\infty \lrate_k^2 < \infty .
% \end{align}

Condition \eqref{eq: summable learning rates outside def} ensures that learning rates eventually remain smaller than 1. Since we study the asymptotic behaviour of \eqref{eq: complete difference inclusion},
the assumption $\alpha_k \in (0,1)$ entails no loss of generality.
So $\alpha_k \update_k(s,a) \le \update_k(s,a)$.
% for every state-action pair, and .
% and condition \eqref{eq: xtn robbins-monro second} implies that $\lrate_k \leq 1$ for sufficiently large $k$.
As a result, condition \eqref{eq: condition infinite visits} is redundant because \eqref{eq: condition infinite effective learning rate} always implies \eqref{eq: condition infinite visits}.
% \eqref{eq: xtn robbins-monro first} guarantees condition \eqref{eq: condition infinite visits}.
% It also guarantees \eqref{eq: condition infinite visits} since
% % , for all state-action pairs,
% \begin{align}
% \label{eq: rm implies infinite visits}
%     \sum_{k \geq 0} \update_k (s,a) 
%     \geq \sum_{k \geq 0} \lrate_k \update_k (s,a) 
%     = \infty.
% \end{align}
The converse is not true.
For instance, if $\lrate_k = \update_k(s,a) = 1/k$ for some state-action pair, the estimate $q_k(s,a)$ may converge prematurely because the product $\prod_{i=k}^{k+n-1} (1-\lrate_i \update_i(s,a))$ would not decay to 0,
despite being updated infinitely often and satisfying condition \eqref{eq: robbins-monro first}.

Overall, \eqref{eq: condition infinite effective learning rate} and \eqref{eq: summable learning rates outside def} are the minimal conditions on the learning rates and update distributions so that $q_{k+n}$ converges to $q_\pi$. We refer to these conditions as the \textit{modified Robbins-Monro conditions}.

\begin{definition}[Modified Robbins-Monro conditions]
\label{asp: stochastic conditions on effective learning rate}
    For every state $s \in \sset$ and action $a \in \aset(s)$,
    the sequences of learning rates $(\lrate_k)_{k \geq 0}$ and update distributions $(\update_k)_{k \geq 0}$ satisfy
        \begin{align}
            % \label{eq: xtn robbins-monro first}
            % \sum_{k\geq0} \lrate_k \update_k(s,a) &= \infty , 
            % \\
            % \label{eq: xtn robbins-monro second}
            % \sum_{k\geq0} \lrate_k^2 \update_k^2(s,a) &< \infty .
            \label{eq: xtn robbins-monro first}
            &\sum_{k = 0}^{\infty} \lrate_k \update_k(s,a) = \infty , 
            \\
            \label{eq: xtn robbins-monro second}
            &\sum_{k = 0}^{\infty} \lrate_k^2 < \infty .
        \end{align}
\end{definition}

\vspace{-0.5cm}

% Further, condition \eqref{eq: xtn robbins-monro second}
% guarantees that for all state-action pairs
% \begin{align}
% \label{eq: rm implies sqared eff lr is bounded}
%     \sum_{k \geq 0} \lrate_k^2 \update_k^2 (s,a) 
%     \leq \sum_{k \geq 0} \lrate_k^2 
%     < \infty .
% \end{align}

The update distributions $(\update_k)$ may depend on the policies $(\pol_k)$ and thus be unknown before running MC-O-PI.
Nevertheless, it is possible to select learning rates, start distributions and update functions so that the modified Robbins-Monro conditions hold.
For initial-visit updates, it suffices to verify that the sequences $(\lrate_k)$ and $(\start_k)$ satisfy these conditions.
If they do, the conditions remain valid under first-visit updates, regardless of the policies $(\pol_k)$.


% \textcolor{red}{how to remove ES from these assumptions? how to choose for RM or modified RM?}

% Altogether, throughout this document, we 

% \begin{assumption}[Standing assumptions]
% \label{asp: standing}
% Throughout this paper, we assume:
% \begin{enumerate}
%     \item \textbf{\sffamily Modified Robbins-Monro learning rates.}
%     % The learning rates $(\lrate_k)$ satisfy
%     \begin{align}
%         0 \le \alpha_k \mu_k(s,a) \le 1 ,
%         \qquad
%         \sum_{k =0}^\infty \alpha_k \mu_k(s,a) = \infty , 
%         \qquad
%         \sum_{k =0}^\infty \alpha_k^2 < \infty .
%     \end{align}
    
%     \item \textbf{\sffamily Unbiased return estimates.} The update functions $(f_k)$ are initial- of first-visit updates.

%     \item \textbf{\sffamily Contractive MDP.}
%     % The Bellman operators $T_\pi$ and $T$ are monotonic and are $\gamma$-contractions in the $L_\infty$-norm.
%     Either the discount factor satisfies $\gamma \in [0, 1)$ or each simulation of the MDP reaches a terminal state almost surely.
% \end{enumerate}
% \end{assumption}


%%%%%%%%%%%%%%%%%%%%%%%%%%%%%%%%%%%%%%%%%
%%%%%%%%%%%%%%%%%%%%%%%%%%%%%%%%%%%%%%%%%
%%%%%%%%%%%%%%%%%%%%%%%%%%%%%%%%%%%%%%%%%
\section{Basic properties}
\label{sec: mcopi behaviour}

% \mytodo{be very clear about when need $(q_k)$ to be generated by MC-O-PI and when consider a random sequence $(q_k)$}

Due to the stochastic episode simulations and asynchronous updates of state-action pairs, the difference inclusion \eqref{eq: complete difference inclusion} that underpins MC-O-PI generates two random sequences: $(q_k)$ and $(\pol_k)$.
We characterise their convergence using the concept of \textit{convergence with probability one}.

% see 
% theorem 4.1.3 in https://doi.org/10.1007/978-3-030-40183-2
% theorem in Tao's Random Matrix theory
% Grimmett - Probability
% Wikipedia
% https://terrytao.wordpress.com/wp-content/uploads/2011/02/matrix-book.pdf

\begin{definition}[Convergence with probability one]
Let $(\xset, d)$ be a metric space.
A sequence $(x_k)_{k \geq 0}$ taking random values in $\xset$ converges with probability one to a point $x \in \xset$ if
\begin{align*}
    \prob \left( \limsup_{k \to \infty} d( x, x_k) \leq \epsilon \right)
    % =
    % \prob \big( 
    % \exists K , \forall k \geq K : \| x - x_k \| \leq \epsilon 
    % \big)
    = 1
\end{align*}
for every $\epsilon>0$.
The same sequence converges with probability one to a set $\xset' \subset \xset$ if
\begin{align*}
    \prob \left( \limsup_{k \to \infty} \inf_{x \in \xset'} d( x, x_k) \leq \epsilon \right)
    % =
    % \prob \big( \exists K , \forall k \geq K : \inf_{x \in \xset} \| x - x_k \| \leq \epsilon \big)
    = 1
\end{align*}
for every $\epsilon>0$.
\end{definition}
\vspace{-0.5cm}

This notion of convergence, also called \textit{almost sure convergence}, means that while events where the sequence $(x_k)$ does not converge to $x$ or $\xset'$ may exist, they collectively have a probability of zero.
This definition allows us to describe some key properties of the sequences $(q_k)$ and $(w_k)$ generated by MC-O-PI in the function space $\qset$ with metric $d(q, q') = \|q-q'\| = \max_{s,a} | q(s,a) - q'(s,a) |$.

% These properties follow from the martingale characteristics of the noise derived in \autoref{thm: properties of noise}.
% In this definition however $w_k$ is inversely proportional to $\update_k$. Thus if for some state-action pair $\update_k(s,a)$ is small, $w_k(s,a)$ may be large.
% This effect is neutralised because $w_k$ is scaled by $\update_k$ at each step of the difference inclusion \eqref{eq: complete difference inclusion}.
% Consequently, the quantities of interest are the bias and more importantly the variance of the product $\update_k w_k$ rather than $w_k$ alone.
% The following lemma shows that the product $\update_k w_k$ has no bias and a bounded variance.

By leveraging the properties of the stochastic perturbations (\autoref{thm: properties of noise}) and the modified Robbins-Monro conditions (\autoref{asp: stochastic conditions on effective learning rate})
we can apply standard martingale and stochastic approximation arguments to establish fundamental properties of the difference inclusion \eqref{eq: complete difference inclusion} that underlies MC-O-PI.
The resulting properties presented below rely extensively on theorems from Chapter 4 of \citep{Bertsekas1996Neuro-DynamicProgramming}.


% \begin{lemma}
\label{eq: properties of noise}
    There exist two constants $c_1, c_2 \in \real_{\geq 0}$ independent of $\update$
    such that
    for every point $q \in \qset$, policy $\pol \in \polset$, state $s \in \sset$ and action $a \in \aset(s)$,
    the noise $w$ defined in equation \eqref{eq: definition combined noise} satisfies
    \begin{align}
        \label{eq: noise prop exp}
        &\E\big[ \update(s,a) w(s,a) \big] = 0 ,
        \\
        \label{eq: noise prop var}
        % &\E\big[ \update^2(s,a) w^2(s,a) \big] \leq c_1 + c_2 \| q \|^2 .
        &\E\big[ w^2(s,a) \big] \leq c_w(1 + \| q \|^2) .
\end{align}
\end{lemma}
\mytodo{update with proof from paper}

\begin{proof}
\textcolor{red}{check if the paper proof does not rely on the paper's standing assumptions}

Fix $k\ge 0$ and $(s,a)$. When $\mu_k(s,a)=0$, the coordinate is not updated. By our convention $w_k(s,a)=0$. So equations \eqref{eq: muw_mds} and \eqref{eq: muw_second_moment} hold trivially. Assume henceforth that $\mu_k(s,a)>0$.

Recall that $u_k=\mu_k+v'_k$.
% with $\mu_k=\E[u_k\mid \mathcal F_k]$.
The definition of $w_k$ in \eqref{eq: definition combined noise} thus yields
\begin{align}
\label{eq: algebra_rewrite_muw}
\mu_k w_k
&= \mu_k v_k + (u_k-\mu_k)(q_{\pi_k}-q_k+v_k) \nonumber\\
&= u_k v_k + (u_k-\mu_k)(q_{\pi_k} - q_k).
\end{align}
Taking conditional expectations and using the facts that $q_k$ and $q_{\pi_k}$ are $\mathcal F_k$-measurable and $\mu_k=\E[u_k\mid \mathcal F_k]$, we get
\begin{align}
\label{eq: refactor noise}
\E[\mu_k w_k \mid \mathcal F_k] &= \E[u_k v_k \mid \mathcal F_k] + \big(q_{\pi_k}-q_k\big)\E[u_k-\mu_k\mid \mathcal F_k] 
\nonumber \\
&= \E[u_k v_k\mid \mathcal F_k] .
\end{align}
% because the second term vanishes from the definition $\mu_k = \E[u_k \mid \mathcal F_k]$.
Using the law of total expectation, we can factor out the indicator $u_k(s,a)$:
\begin{align}
\label{eq: showter expectation}
\E[u_k(s,a)v_k(s,a)\mid \mathcal F_k] 
= \E\big[
u_k(s,a)\,\E[v_k(s,a)\mid \mathcal F_k,\ u_k(s,a)]\mid \mathcal F_k
\big] .
\end{align}
If $(s,a)$ is not updated after the $k^{\text{th}}$ episode we have $u_k(s,a)=0$, and arbitrarily define $v_k(s,a) = 0$ to maintain consistency with $w_k(s,a)=0$.
Therefore, 
equation \eqref{eq: muw_mds} is equivalent to $\E[v_k(s,a)\mid \mathcal F_k, u_k(s,a)=1] = 0$.

With initial- or first-visit updates, the event that a state-action pair is updated, i.e. $u_k(s,a) = 1$, depends entirely on the episode's history prior to that first occurrence of $(s,a)$. By the Markov property, the simulated episode from that point onward depends exclusively on the pair $(s,a)$ and the current policy $\pi_k$. In other words, it is completely independent of the events that led to the first visit of $(s,a)$. Consequently, the return noise $v_k(s,a)$ behaves exactly as the intrinsic zero-mean noise of a fresh rollout from $(s,a)$, guaranteeing that $\E[v_k(s,a) \mid \mathcal{F}_k, u_k(s,a)=1] = 0$.
As a result, 
% equations \eqref{eq: refactor noise} and \eqref{eq: showter expectation} indicate that 
the unbiased noise property holds for every state-action pair at every time step.


Applying the standard inequality $(x+y)^2 \le 2x^2 + 2y^2$ to equation \eqref{eq: algebra_rewrite_muw}
and noting that $u_k^2 \le 1$ and $(u_k-\mu_k)^2 \le 1$, we find
\[
(\mu_k w_k)^2 \le 2v_k^2 + 2(q_{\pi_k}-q_k)^2 .
\]
Taking the conditional expectation yields
\[
\E\!\left[\mu_k^2 w_k^2 \mid \mathcal F_k\right] \le 2\E\!\left[v_k^2 \mid \mathcal F_k\right] + 2(q_{\pi_k}-q_k)^2
\]
Since $v_k$ is the noise on the observed returns after simulating an episode of the MDP with policy $\pi_k$,
there exists a constant $C_v<\infty$ such that $\E[v_k^2(s,a) \mid \mathcal F_k] \le C_v$ for every state-action pair at every time step. 
Moreover, the action value $q_{\pi_k}(s,a)$ is uniformly bounded by the constant $C_\pi \coloneqq \max_{s,a} r(s,a)/(1-\discount)$. Applying $(x-y)^2 \le 2x^2 + 2y^2$, we see that
\[
\big( q_{\pi_k}(s,a) - q_k(s,a) \big)^2 \le 2 q_{\pi_k}(s,a)^2 + 2 q_k(s,a)^2 \le 2C_\pi^2 + 2 \max_{s,a} |q_k(s,a)|^2
\]
% where $\|\cdot\|$ is any norm dominating the coordinate magnitude.
Substituting these into our expectation bound yields
\[
\E [ \mu_k^2(s,a) w_k^2(s,a) \mid \mathcal F_k ] \le 2C_v + 4C_\pi^2 + 4\|q_k\|_\infty^2
\]
Setting $C \coloneqq \max \{ 2C_v + 4C_\pi^2, 4\}$ gives \eqref{eq: muw_second_moment},
independently of $\mu_k$.


% \hline
\textcolor{red}{false}
Since the variables $w_1$ and $w_2$ defined in \eqref{eq: def noise qpol} and \eqref{eq: def noise u} are independent, we know that
\begin{align*}
    \E \big[ w_1 w_2 \big]
    =     \E \big[ w_1 \big]  \E \big[ w_2 \big] = \zerofun 
\end{align*}
where $\zerofun$ is the constant function $\zerofun(s,a) = 0$.
Thus it follows that
\begin{align*}
    \E \big[ \update w \big] 
    &= \E \big[ \update w_1 + w_2 ( q_\pol - q + w_1 ) \big] 
    \\
    &= \update \E[w_1] + \E [w_2] \big( q_\pol - q + \E[ w_1 ] \big) 
    \\
    &= \zerofun.
\end{align*}
Further, we know that for all state-action pairs $0 \leq \update(s,a) \leq 1$ and $-1 \leq w_2(s,a) \leq 1$. Thus we obtain that
\begin{align*}
    \E \big[ \update^2 w^2 \big]
    &= \E \big[ ( \update w_1 + w_2 ( q_\pol - q + w_1 ) )^2 \big]
    \\
    % &= \E \Big[ w_1^2 + \frac{w_2^2}{\update^2} ( q_\pol - q + w_1 )^2 + 2 \frac{w_1 w_2}{\update} ( q_\pol - q + w_1 ) \Big]
    % \\
    % &= \E \big[ w_1^2 \big] + \E \Big[ \frac{w_2^2}{\update^2} ( q_\pol - q + w_1 )^2 \Big] + 2 \E \Big[ \frac{w_1 w_2}{\update} ( q_\pol - q ) \Big] + 2 \E \Big[ \frac{w_1^2 w_2}{\update} \Big]
    % \\
    &= \update^2 \E \big[ w_1^2 \big] + \E \big[ w_2^2 ( q_\pol - q + w_1 )^2 \big] +  2 \E \big[ \update w_1 w_2 ( q_\pol - q + w_1 ) \big] 
    \\
    &= \update^2 \E \big[ w_1^2 \big] + \E \big[ w_2^2 ( q_\pol - q + w_1 )^2 \big]
    \\
    &\leq \E \big[ w_1^2 \big] + \E \big[ ( q_\pol - q + w_1 )^2 \big]
    \\
    &= \E \big[ w_1^2 \big] + (q_\pol - q)^2 + \E \big[ w_1^2 \big] + 2 (q_\pol - q) \E \big[ w_1 \big]
    \\
    &\leq c_1 \onefun + c_2 q^2
    % \\
    % &\leq (c_1 + c_2 \|q\|^2) \onefun
\end{align*}
for some constants $c_1, c_2 \in \real_{\geq 0}$ when every action-value function $q_\pol$ is bounded.
Finally, \eqref{eq: noise prop var} follows from the inequality $q(s,a)^2 \leq \|q\|^2$.
\end{proof}


%%%%%%%%%%%%%%%%%%%%%%%%%%%%%%%%%%%%%%%%%
%%%%%%%%%%%%%%%%%%%%%%%%%%%%%%%%%%%%%%%%%
% \subsection{Solutions are bounded}

% \begin{lemma}
% \label{thm: bounded mcopi solutions}
% If the effective learning rates $(\lrate_k \update_k)$ satisfy the Robbins-Monro conditions (\autoref{asp: stochastic conditions on effective learning rate})
% and solutions of MC-O-PI only visit a subset of policies $\polset' \subseteq \polset$,
% then solutions are almost surely bounded, meaning that
% \begin{align}
%     &\limsup_{k \to \infty} q_k(s,a) \leq \max_{\pol \in \polset'} q_{\pol}(s,a)
%     \\
%     &\liminf_{k \to \infty} q_k(s,a) \geq \min_{\pol \in \polset'} q_{\pol}(s,a)
% \end{align}
% \end{lemma}
% \begin{proof}
%     just take proof of other lemma
% \end{proof}


% Given that the difference inclusion \eqref{eq: complete difference inclusion} is essentially a stable linear system with a variable set point that is perturbed by noise, and that the eigenvalues of the system are real
% The following lemma shows that solutions of MC-O-PI converge to $\qset_\infty$.

% \begin{lemma}
% \label{thm: bounded solution}
%     Solutions of MC-O-PI, as defined in \eqref{eq: complete difference inclusion}, converge to the region $\qset_\infty$
%     % defined in \eqref{eq: long term region} 
%     with probability one
%     if the learning rates $(\lrate_k)_{k \geq 0}$ and update distributions $(\update_k)_{k \geq 0}$ satisfy the modified Robbins-Monro conditions (\autoref{asp: stochastic conditions on effective learning rate}).
%     % For any $\epsilon>0$ there exists a time step $K$ when the sequence $(q_\kpz)$ generated by MC-OPI is contained for all $k \geq K$ inside the region
% \end{lemma}
\begin{lemma}
    [Bounded solutions]
    \label{thm: bounded limit sets}
    Let $\polset_\infty$ denote the set of policies that appear infinitely often in a sequence $(\pi_k)_{k \ge 0}$ generated by MC-O-PI (\autoref{def: mcopi}).
    % \begin{align}
    %     \polset_\infty \coloneqq \bigcap_{t = 0}^{\infty} \text{cl} \{\pi_k : k \geq t \} ,
    % \end{align}
    If the learning rates $(\lrate_k)_{k \geq 0}$ and update distributions $(\update_k)_{k \geq 0}$ satisfy the modified Robbins-Monro conditions (\autoref{asp: stochastic conditions on effective learning rate}),
    % Under \autoref{asp: standing} 
    the corresponding sequence $(q_k)_{k \ge 0}$ converges almost surely to the set
    \begin{align*}
        % \qset_\infty 
        % \coloneqq
        \Big\{
        q \in \qset : 
        \min_{\pi \in \polset_\infty} q_\pi(s,a) \leq q(s,a) \leq \max_{\pi \in \polset_\infty} q_\pi(s,a),
        % \forall s \in \sset, \forall a \in \aset
        \forall \, (s,a)
        \Big\} .
    \end{align*}
    Consequently,
    % for every solution $(q_k)$ of MC-O-PI,
    there exists an almost surely finite $c_Q \in [0, \infty)$ such that for all $k \geq 0$
        \begin{align*}
            \| q_k \|_\infty^2 < c_Q .
        \end{align*}
\end{lemma}
\vspace{-0.6cm}
% \begin{proof}
%     Let $(q_k)$ and $(\pol_k)$ be solutions of the difference inclusion \eqref{eq: complete difference inclusion} with initial condition $q_0$, learning rates $(\lrate_k)$, update distributions $(\update_k)$ and noise terms $(w_k)$.
%     At every step, the estimate $q_k$ is updated towards action values $q_{\pol_k}$ that satisfy $q_{\worst{\pol}} \leq q_{\pol_k} \leq q_{\best{\pol}}$.
%     % \begin{align*}
%     %     q_{\worst{\pol}} \leq q_{\pol_k} \leq q_{\best{\pol}} .
%     % \end{align*}
%     Therefore, define the sequences $(x_k)$ and $(y_k)$ as
%     \begin{align}
%         x_\kpo &= x_k + \lrate_k \update_k (q_{\worst{\pol}} - x_k + w_k) ,
%         \\
%         y_\kpo &= y_k + \lrate_k \update_k (q_{\best{\pol}} - y_k + w_k) ,
%     \end{align}
%     and initialised at $x_0 = y_0 = q_0$. At every step we have
%     \begin{align*}
%         x_k \leq q_k \leq y_k .
%     \end{align*}
%     Proposition~4.1 and Example~4.3 in \citep{Bertsekas1996Neuro-DynamicProgramming} show that the sequences $(x_k)$ and $(y_k)$ converge with probability one to $q_{\worst{\pol}}$ and $q_{\best{\pol}}$, respectively, because the noise satisfies conditions \eqref{eq: noise prop exp} and \eqref{eq: noise prop var}.
%     As a result, we conclude that almost surely
%     \begin{align*}
%         q_{\worst{\pol}} = \liminf_{k \to \infty} q_k 
%         \leq q_k 
%         \leq \limsup_{k \to \infty} q_k = q_{\best{\pol}} ,
%     \end{align*}
%     % As a result, for every $\epsilon>0$ there exists a time step $K$ such that 
%     % \begin{align*}
%     %     q_{\worst{\pol}} - \epsilon \onefun \leq x_k \leq q_k \leq y_k \leq q_{\best{\pol}} + \epsilon \onefun
%     % \end{align*}
%     % for all $k \geq K$ with probability one.
%     which implies that solutions converge to $\qset_\infty$ with probability one.
% \end{proof}
\begin{proof}
    Let $(q_k)_{k \ge 0}$ and $(\pol_k)_{k \ge 0}$ be solutions of the difference inclusion in \eqref{eq: complete difference inclusion} with initial condition $q_0$, learning rates $(\lrate_k)_{k \ge 0}$, update distributions $(\update_k)_{k \ge 0}$ and stochastic perturbations $(w_k)_{k \ge 0}$.
    
    By assumption, there exists a finite time $K$ such that
    the sequence $(\pol_k)_{k \ge K}$ only contains policies in $\polset_\infty$.
    % the estimate $q_k$ is updated towards action values $q_{\pi_k}$ where $\pi_k \in \polset_\infty$.
    Thus for all $k \ge K$, each estimate $q_k(s,a)$ is updated towards a value $q_{\pol_k}(s,a)$ that satisfies
    \begin{align}
        \min_{\pi \in \polset_\infty} q_\pi(s,a) \leq q_{\pi_k}(s,a) \leq \max_{\pi \in \polset_\infty} q_\pi(s,a) .
    \end{align}
    
    Therefore, define the scalar sequences $(x_k)_{k \ge K}$ and $(y_k)_{k \ge K}$ as
    \begin{align}
        x_\kpo &= x_k + \lrate_k \update_k(s,a) 
        \big( \min_{\pi \in \polset_\infty} q_\pi(s,a) - x_k + w_k(s,a) \big) ,
        \\
        y_\kpo &= y_k + \lrate_k \update_k(s,a) \big( \max_{\pi \in \polset_\infty} q_\pi(s,a) - y_k + w_k(s,a) \big) ,
    \end{align}
    and initialised at $x_K = y_K = q_K(s,a)$. From time step $K$ onwards we have
    \begin{align*}
        x_k \leq q_k(s,a) \leq y_k .
    \end{align*}
    Proposition~4.1 and Example~4.3 in \citep{Bertsekas1996Neuro-DynamicProgramming} show that $x_k$ and $y_k$ converge almost surely to $\min_{\pi \in \polset_\infty} q_\pi(s,a)$ and $\max_{\pi \in \polset_\infty} q_\pi(s,a)$, respectively, because the noise satisfies conditions \eqref{eq: muw_mds} and \eqref{eq: muw_second_moment}.
    As a result, we conclude that almost surely
    \begin{align*}
        &\liminf_{k \to \infty} q_k(s,a)
        \ge \min_{\pi \in \polset_\infty} q_\pi(s,a)
        % \leq q_k(s,a) 
        \\
        &\limsup_{k \to \infty} q_k(s,a)
        \le \max_{\pi \in \polset_\infty} q_\pi(s,a)
    \end{align*}
    for every state-action pair,
    which proves the first part of the lemma.

    Further,
    given that there are a finite number of deterministic policies, $q_k$ converges almost surely to a bounded set.
    As a result, 
    % that sequence is almost surely bounded, meaning that 
    its norm over all $k \geq 0$ is almost surely finite. So there exists an almost surely finite constant $c_Q$ that is strictly greater than $\sup_{k \geq 0} \|q_k\|_\infty^2$.
\end{proof}

\vspace{-0.4cm}

\autoref{thm: bounded limit sets} implies that MC-O-PI solutions converge almost surely to the long-term region $\qset_\infty$ defined as
\begin{align}
\label{eq: long term region}
    \qset_\infty \coloneqq \big\{ q \in \qset : q_{\worst{\pol}} \leq q \leq q_{\best{\pol}} \big\} ,
\end{align}
where $\worst{\pol}$ and $\best{\pol}$ are the worst and best policies of the MDP.
% and the inequalities hold elementwise.
Therefore, when studying the asymptotic behaviour of MC-O-PI, it is sufficient to consider solutions that start in the long-term region $\qset_\infty$.

% \begin{corollary}
\label{thm: solution of mcopi is bounded}
    If the learning rates $(\lrate_k)_{k \geq 0}$ and update distributions $(\update_k)_{k \geq 0}$ satisfy the modified Robbins-Monro conditions (\autoref{asp: stochastic conditions on effective learning rate}),
    then for every solution $(q_k)_{k \geq 0}$ of MC-O-PI, as defined in \eqref{eq: complete difference inclusion}, there exists almost surely a constant $c \in \real_{\geq 0}$ such that
    \begin{align}
        \| q_k \|^2 < c
    \end{align}
    for all $k \geq 0$.
\end{corollary}
\begin{proof}
    Consider any bounded initial condition $q_0 \in \qset$.
    By assumption, the action values of all policies are bounded.
    Furthermore, \autoref{eq: properties of noise} establishes that the variance of $\update_k w_k$ is bounded.
    Therefore, if $q_k$ is bounded, the update $\lrate_\kpz \update_k ( q_{\pol_k} - q_k + w_k )$ is bounded, and consequently $q_\kpo$ is bounded.
    \autoref{thm: bounded solution} then shows that the sequence $(q_k)$ does not grow indefinitely as it establishes that for every $\epsilon>0$, there exists almost surely a time step $K$ such that
    \begin{align*}
        q_{\worst{\pol}} - \epsilon \onefun \leq q_k \leq q_{\best{\pol}} + \epsilon \onefun 
    \end{align*}
    for all $k \geq K$.
    As a result, the solution is bounded from zero to $K$ and from $K$ to infinity with probability one.
    Thus, there exists almost surely a constant $c \in \real_{\geq 0}$ such that $\| q_k \|^2 < c$
    for all $k \geq 0$.
\end{proof}


%%%%%%%%%%%%%%%%%%%%%%%%%%%%%%%%%%%%%%%%%
%%%%%%%%%%%%%%%%%%%%%%%%%%%%%%%%%%%%%%%%%
% \subsection{Effects of noise become arbitrarily small}

Having established that solutions converge to $\qset_\infty$, we now show that the effects of the perturbations are asymptotically negligible, meaning they cannot prevent solutions from converging to a point in $\qset_\infty$.
% This follows from the next lemma that proves that the perturbations $(\lrate_k \update_k w_k)$ converge to zero.

% \begin{lemma}
% \label{thm: effects of noise decrease}
%     For every initial condition $q_0 \in \qset$, every possible sequence $(\lrate_k \update_k w_k)_{k \geq 0}$ generated by MC-O-PI, as defined in \eqref{eq: complete difference inclusion}, converges to zero with probability one
%     if the learning rates $(\lrate_k)_{k \geq 0}$ and update distributions $(\update_k)_{k \geq 0}$ satisfy the modified Robbins-Monro conditions (\autoref{asp: stochastic conditions on effective learning rate}).
%     % meaning that, for every $\epsilon>0$,
%     % \begin{align}
%     %     \prob \big( \limsup_{k \to \infty} \| \lrate_k \update_k w_k \| \leq \epsilon \big) = 1 .
%     % \end{align}

%     For every initial condition $q_0 \in \qset$, 
%     the series $\sum_{k=0}^\infty \lrate_k \update_k w_k$ generated by MC-O-PI converges almost surely to a finite limit under \autoref{asp: standing}.
%     Hence the sequence $(\lrate_k \update_k w_k)_{k \ge 0}$ converges almost surely to zero.
% \end{lemma}
\begin{lemma}
[Bounded total perturbations]
\label{thm: effects of noise decrease}
    For every initial condition $q_0 \in \qset$,
    every state $s \in \sset$ 
    and every action $a \in \aset(s)$,
    if the learning rates $(\lrate_k)_{k \geq 0}$ satisfy 
    $\sum_{k = 0}^{\infty} \alpha_k^2 < \infty$,
    then the sum $\sum_{k=0}^\infty \lrate_k w_k(s,a)$ generated by MC-O-PI (\autoref{def: mcopi}) is almost surely finite.
    Hence the sequence $(\lrate_k w_k(s,a))_{k \ge 0}$ converges almost surely to zero.
\end{lemma}

% \begin{proof}
% % Similarly to the proof of Lemma~4.1 in \cite{Bertsekas1996Neuro-DynamicProgramming}
% Let $(q_k)$ be a solution of the difference inclusion \eqref{eq: complete difference inclusion} with initial condition $q_0$, learning rates $(\lrate_k)$, update distributions $(\update_k)$ and noise terms $(w_k)$.
% \autoref{eq: properties of noise} shows that there exist two positive constants $c_1$ and $c_2$ such that, for all state-action pairs and $k \geq 0$,
% \begin{align*}
%     &\E \big[ \update_k(s,a) w_k(s,a) \big] = 0 ,
%     \\
%     &\E \big[ \update_k^2(s,a) w_k^2(s,a) \big] \leq c_1 + c_2 \|q_k\|^2 ,
% \end{align*}
% while \autoref{thm: solution of mcopi is bounded} establishes that there exists almost surely a positive constant $c_3$ such that, for all $k \geq 0$,
% \begin{align*}
%     \|q_k\|^2 < c_3 .
% \end{align*}

% Now, define a sequence $(x_k)$ as 
% \begin{align}
%     x_k = \sum_{k' = 0}^{k} \lrate_{k'} \update_{k'} w_{k'} .
%     % x_\kpo = x_k + \lrate_k \update_k w_k
% \end{align}
% We observe that for all $k \geq 0$
% \begin{align}
%     \notag
%     \E \big[ x_k \mid x_0, \dots, x_{k-1} \big] 
%     &= \E \big[ x_{k-1} + \lrate_k \update_k w_k \mid x_0, \dots, x_{k-1} \big] 
%     \\
%     \notag
%     &= x_{k-1} + \lrate_k \E \big[ \update_k w_k \mid x_0, \dots, x_{k-1} \big] 
%     \\
%     \label{eq: first for martingale}
%     &= x_{k-1} .
% \end{align}
% It also almost surely holds that
% \begin{align}
%     \notag
%      \E \big[ x_k^2 \mid x_{k-1} \big]
%      &= \E \big[ ( x_{k-1} + \lrate_k \update_k w_k )^2 \mid x_{k-1} \big]
%      \\
%      \notag
%      &= x_{k-1}^2 
%      + \lrate_k^2 \E \big[ \update_k^2 w_k^2 \mid x_{k-1} \big]
%      + 2 \lrate_k x_{k-1} \E \big[ \update_k w_k \mid x_{k-1} \big]
%      \\
%      \notag
%      &\leq x_{k-1}^2 + \lrate_k^2 ( c_1 + c_2 \|q_k\|^2 ) \onefun
%      \\
%      &\leq x_{k-1}^2 + \lrate_k^2 ( c_1 + c_2 c_3 ) \onefun .
% \end{align}
% Since the absolute value satisfies $|x| \leq \onefun + x^2$,
% it almost surely follows that
% \begin{align}   
%     \notag
%     \E \big[ | x_k | \big]
%     &=
%     \E \Big[ \E \big[ | x_k | \mid x_{k-1} \big] \Big]
%     \\
%     \notag
%     &\leq \E \Big[ \onefun + \E \big[ x_k^2 \mid x_{k-1} \big] \Big]
%     \\
%     \notag
%     % &= \E \Big[ \onefun + \E \big[ ( x_{k-1} + \lrate_k \update_k W_k )^2 \mid x_{k-1} \big] \Big]
%     % \\
%     % &= \E \Big[ \onefun +  \Big]
%     % \\
%     % &\leq \E \Big[ \onefun + x_{k-1}^2 + \lrate_k^2 \update_k^2 ( c_2 \onefun + c_3 q_k^2 ) \Big]
%     % \\
%     % &= \onefun + \E \big[ x_{k-1}^2 \big] + \lrate_k^2 \update_k^2 ( c_2 \onefun + c_3 q_k^2 )
%     % \\
%     &\leq \onefun + \E \big[ x_{k-1}^2 + \lrate_k^2 ( c_1 + c_2 c_3 ) \onefun \big]
%     \\
%     \notag
%     &= \onefun + \E \big[ x_{k-1}^2 \big] + \lrate_k^2 ( c_1 + c_2 c_3 ) \onefun
%     \\
%     \notag
%     &= \onefun + \E \Big[ \E \big[ x_{k-1}^2 \mid x_{k-2} \big] \Big] + \lrate_k^2 ( c_1 + c_2 c_3 ) \onefun
%     \\
%     \notag
%     &\leq \cdots
%     \\
%     \notag
%     &\leq \onefun + ( c_1 + c_2 c_3 ) \sum_{k = 0}^{\infty} \lrate_k^2 \onefun
%     \\
%     \label{eq: second for martingale}
%     &< \infty \onefun .
% \end{align}
% From equations \eqref{eq: first for martingale} and \eqref{eq: second for martingale}, the Martingale Convergence Theorem ensures the convergence of the sequence $(x_k)$.
% Since $(x_k)$ contains the partial sums of the series $\sum_{k \geq 0} \lrate_k \update_k w_k$, Corollary~7.2.6 in \citep{Tao2016AnalysisI} implies that the sequence $(\lrate_k \update_k w_k)$ converges to zero with probability one.
% \end{proof}

\begin{proof}
    % Similarly to the proof of Lemma~4.1 in \cite{Bertsekas1996Neuro-DynamicProgramming}
    Let $(q_k)_{k \ge 0}$ be a solution of the difference inclusion in \eqref{eq: complete difference inclusion} with initial condition $q_0$, learning rates $(\lrate_k)_{k \ge 0}$, update distributions $(\update_k)_{k \ge 0}$ and stochastic perturbations $(w_k)_{k \ge 0}$.
    
    % \autoref{eq: properties of noise} shows that there exist two positive constants $c_1$ and $c_2$ such that, for every state-action pair and $k \ge 0$,
    % \begin{align*}
    %     &\E \big[ \update_k(s,a) w_k(s,a) \big] = 0 ,
    %     \\
    %     &\E \big[ \update_k^2(s,a) w_k^2(s,a) \big] \leq c_1 + c_2 \|q_k\|^2 ,
    % \end{align*}
    % while \autoref{thm: bounded solutions} establishes that there exists almost surely a positive constant $c_3$ such that, for all $k \geq 0$,
    % \begin{align*}
    %     \|q_k\|^2 < c_3 .
    % \end{align*}

    Fix a state $s \in \sset$ and an action $a \in \aset(s)$.
    Define a sequence of partial sums $(x_k)_{k \ge 0}$ as 
    \begin{align}
        x_k = \sum_{i = 0}^{k} \lrate_{i} w_{i}(s,a) .
        % x_\kpo = x_k + \lrate_k \update_k w_k
    \end{align}
    % Let $\mathcal{F}_k$ be the filtration defined in \eqref{eq: filtration_def}. 
    By \autoref{thm: properties of noise}, the perturbations are conditionally unbiased, meaning that $\mathbb{E}[\lrate_k w_k(s,a) \mid \mathcal{F}_k] = 0$. Therefore, $(x_k)_{k \ge 0}$ is a martingale adapted to the filtration $(\mathcal{F}_{k})_{k \ge 0}$.
    
    \autoref{thm: properties of noise} and \autoref{thm: bounded limit sets} 
    also imply that there exist almost surely finite constants $c_w, c_Q \in [0, \infty)$ such that the variance of the martingale increments satisfy
    \begin{align*}
    \E \left[(x_k - x_{k-1})^2 \mid \mathcal{F}_k \right] 
    =
    \alpha_k^2 \E [w_k^2(s,a) \mid \mathcal{F}_k] 
    % \le \lrate_k^2 C(1 + \|q_k\|^2)
    \le \lrate_k^2 c_w(1 + c_Q) .
    \end{align*}
    Thus if the learning rates satisfy $\sum_{k=0}^\infty \alpha_k^2 < \infty$, the sum of the variances is finite:
    \[
    \sum_{k=1}^\infty \mathbb{E}[(x_k - x_{k-1})^2] \le c_w(1 + c_Q) \sum_{k=1}^\infty \alpha_k^2 < \infty .
    \]
    Altogether, $(x_k)_{k \ge 0}$ is a martingale with bounded $L^2$ variations. Thus Doob's Martingale Convergence Theorem indicates that $x_k$ converges almost surely to a finite limit. Furthermore, since the sequence of partial sums converges, the individual terms $(\alpha_k w_k(s,a))_{k \ge 0}$ almost surely converge to zero.
\end{proof}

% \autoref{thm: effects of noise decrease} demonstrates that the noise terms cannot counteract the decreasing learning rates and cause $(q_k)$ to be a random walk through $\qset$.


%%%%%%%%%%%%%%%%%%%%%%%%%%%%%%%%%%%%%%%%%
%%%%%%%%%%%%%%%%%%%%%%%%%%%%%%%%%%%%%%%%%
% \subsection{Steps become arbitrarily small}

Since solutions converge to the bounded region $\qset_\infty$ and the effects of noise become arbitrarily small, the updates of the differential inclusion \eqref{eq: complete difference inclusion} become arbitrarily small. In other words, the distance between consecutive elements of a solution $(q_k)$ tends to zero.

\begin{lemma}[Steps become arbitrarily small]
\label{thm: steps become arbitrarily small}
    For every solution $(q_k)_{k \ge 0}$ of MC-O-PI (\autoref{def: mcopi}),
    the sequence $( \|q_\kpo - q_k\| )_{k \geq 0}$ converges to zero with probability one
    if the learning rates $(\lrate_k)_{k \geq 0}$ and update distributions $(\update_k)_{k \geq 0}$ satisfy the modified Robbins-Monro conditions (\autoref{asp: stochastic conditions on effective learning rate}).
    % meaning that, for every $\epsilon>0$,
    % \begin{align}
    %     \prob \big( \limsup_{k \to \infty} \| q_\kpo - q_k \| \leq \epsilon \big) = 1 .
    % \end{align}
\end{lemma}
\vspace{-0.7cm}
% \begin{proof}
% We know from expression \eqref{eq: rm implies sqared eff lr is bounded}, that the effective learning rates $(\lrate_k \update_k)$ converge to zero for all state-action pairs if the modified Robbins-Monro conditions hold.
% \autoref{thm: bounded solution} then establishes that solutions of MC-O-PI converge to the region $\qset_\infty$ with probability one.
% Therefore, for every $\epsilon > 0$, there exists almost surely a time step $K'$ such that, for all $k \geq K'$,
% \begin{align*}
%  \|\lrate_k \update_k (q_{\pol_k} - q_k)\| 
%  &\leq \|\lrate_k \update_k \| \|q_{\pol_k} - q_k\| 
%  \\
%  &\leq \|\lrate_k \update_k \| \|\qopt - q_{\worst{\pol}}\|
%  \\
%  &\leq \frac{\epsilon}{2} .
% \end{align*}
% Further, \autoref{thm: effects of noise decrease} demonstrates that there exists almost surely a time step $K''$ such that, for all $k \geq K''$,
% \begin{align*}
%  \|\lrate_k \update_k w_k \| \leq \frac{\epsilon}{2} .
% \end{align*}
% As a result, for every $\epsilon > 0$,
% there exists almost surely a time step $K \geq \max[K', K'']$ such that, for all $k \geq K$,
% \begin{align*}
%     \| q_\kpo - q_k \| 
%     &= \| \lrate_k \update_k (q_{\pol_k} - q_k + w_k) \|
%     \\
%     &\leq \| \lrate_k \update_k (q_{\pol_k} - q_k) \| + \| \lrate_k \update_k w_k \|
%     \\
%     &\leq \epsilon ,
% \end{align*}
% which indicates that the sequence $(\|q_\kpo - q_k\|)_{k \geq 0}$ converges to zero with probability one.
% \end{proof}
\begin{proof}
    % Since $\update_k(s,a) \le 1$ for every S
    The modified Robbins-Monro conditions and the boundedness of the update distributions ensure that the sequence $(\lrate_k \update_k)_{k \ge 0}$ converges to zero for all state-action pairs.
    Additionally, \autoref{thm: bounded limit sets} implies that solutions of MC-O-PI converge almost surely to the region
    \begin{align}
        \big\{q \in \qset : q_{\worst{\pi}} \le q \le \qopt \big\}
    \end{align}
    where $q_{\worst{\pi}}$ and $\qopt$ are the action values of the worst and the best policies, respectively.
    Therefore, for every $\epsilon > 0$, there exists an almost surely finite time $K'$ such that, for all $k \geq K'$,
    \begin{align*}
     \|\lrate_k \update_k (q_{\pol_k} - q_k)\| 
     &\leq \|\lrate_k \update_k \| \|q_{\pol_k} - q_k\| 
     \\
     &\leq \|\lrate_k \update_k \| \|\qopt - q_{\worst{\pol}}\|
     \\
     &\leq \frac{\epsilon}{2} .
    \end{align*}
    Further, \autoref{thm: effects of noise decrease} demonstrates that there exists an almost surely finite time $K''$ such that, for all $k \geq K''$,
    \begin{align*}
     \|\lrate_k w_k \| \leq \frac{\epsilon}{2} .
    \end{align*}
    As a result, for every $\epsilon > 0$,
    there exists an almost surely finite time $K \geq \max\{K', K''\}$ such that, for all $k \geq K$,
    \begin{align*}
        \| q_\kpo - q_k \| 
        &= \| \lrate_k ( \update_k (q_{\pol_k} - q_k ) + w_k) \|
        \\
        &\leq \| \lrate_k \update_k (q_{\pol_k} - q_k) \| + \| \lrate_k w_k \|
        \\
        &\leq \epsilon ,
    \end{align*}
    indicating that the sequence $(\|q_\kpo - q_k\|)_{k \ge 0}$ converges almost surely to zero.
\end{proof}

%%%%%%%%%%%%%%%%%%%%%%%%%%%%%%%%%%%%%%%%%
% %%%%%%%%%%%%%%%%%%%%%%%%%%%%%%%%%%%%%%%%%
% \subsubsection{Suboptimal fixed points only appear on the boundary}

% The Bellman optimality principle indicates that only the action values of the optimal policy $\best{\pol}$ lie in their own region $\region(\best{\pol})$.

% As a result, \eqref{eq: complete difference inclusion}

% MC-O-PI cannot converge to a suboptimal policy

% \begin{lemma}
\label{thm: fixed points are on the boundary}
    If the learning rates $(\lrate_k)$ and the update maps $(\updateset_k)$ satisfy \autoref{asp: stochastic conditions on effective learning rate}, 
    a sequence $(q_k)$ generated by MC-O-PI converges with probability one to a suboptimal fixed point 
    if and only if that point lies on the boundary between several policy regions.
\end{lemma}
\begin{proof}
    Let the sequences $(q_k)$ and $(\pol_k)$ be generated with MC-O-PI, 
    and suppose $(q_k)$ converges with probability one to a suboptimal fixed point $\hat{q}$. 
    If $\hat{q}$ is not on the boundary between several regions, there exists a unique policy $\hat{\pol}$ and an $\epsilon > 0$ such that 
    \begin{align*}
        \forall q \in \qset : 
        \quad
        \|\hat{q} - q\| \leq \epsilon
        \implies
        \avalset(q) = \{q_{\hat{\pol}}\} .
        % \forall s \in \sset
        % \hat{\pol} = \argmax_{\pol \in \polset} \hat{q}(s, \pol) = \argmax_{\pol \in \polset} q(s, \pol) .
    \end{align*}
    Furthermore, the definition of convergence implies a time step $K$ such that for all $k \geq K$
    \begin{align*}
        \|\hat{q} - q_k\| \leq \epsilon
    \end{align*}
    with probability one.
    As a result, for all $k \geq K$ 
    % it also holds that $q_{\pol_k} = q_{\hat{\pol}}$ with probability one.
    system \eqref{eq: complete difference inclusion} simplifies to
    \begin{align}
        \label{eq: mcopi for polk}
        q_\kpo \in \big\{ q_k + \lrate_k \update_k (q_{\hat{\pol}} - q_k + w_k) : \update_k \in \updateset_k(q_k) \big\} .
    \end{align}
    Proposition~4.1 and Example~4.3 in \citep{Bertsekas1996Neuro-DynamicProgramming}  guarantee that, under \autoref{asp: stochastic conditions on effective learning rate}, every sequence generated by system \eqref{eq: mcopi for polk} converges to $q_{\hat{\pol}}$ with probability one, which implies that $\hat{q} = q_{\hat{\pol}}$.
    However, as noted in \autoref{rmk: only qopt is inside its own region}, if $\hat{q}$, and thus $q_{\hat{\pol}}$, are suboptimal, then $\hat{q} \notin \region(\hat{\pol})$, which contradicts the initial assumption.
    As a result, if the sequence $(q_k)$ converges with probability one to a suboptimal fixed point, it must lie on the boundary between several regions.
\end{proof}


% \autoref{thm: fixed points are on the boundary} indicates that ...
% So if can check that several policies are greedy on $q$, then might not have converged to optimal solution.
% Conversely, if on the long term a single policy is greedy on $q$, then are guaranteed that have found optimal policy.
% We will now present the additional conditions where previous research has shown that every solution of MC-O-PI converges to $q_{\best{\pol}}$ and $\best{\pol}$.


%%%%%%%%%%%%%%%%%%%%%%%%%%%%%%%%%%%%%%%%%
%%%%%%%%%%%%%%%%%%%%%%%%%%%%%%%%%%%%%%%%%
%%%%%%%%%%%%%%%%%%%%%%%%%%%%%%%%%%%%%%%%%
\section{Convergence}
\label{sec: existing convergence proofs}


\autoref{thm: bounded limit sets}, \autoref{thm: effects of noise decrease} and \autoref{thm: steps become arbitrarily small} establish that when MC-O-PI satisfies the modified Robbins-Monro conditions, solutions remain bounded, perturbations cannot prevent convergence and updates become arbitrarily small.
Despite these strong properties, \cite{Bertsekas1996Neuro-DynamicProgramming} have constructed a counterexample where MC-O-PI converges to a suboptimal fixed point when only the modified Robbins-Monro conditions hold.
Thus additional statistical regularity is necessary to guarantee that MC-O-PI converges to $\qopt$.

Prior work has addressed this issue by imposing conditions on either the parameters of MC-O-PI or on the structure of the MDP.
\citet{Tsitsiklis2002OnIteration}, \citet{Chen2018OnProblem} and \citet{Liu2021OnStarts} demonstrate convergence to $\qopt$ when the update distributions are uniform over all state-action pairs,
% \cite{Lubars2021OptimisticStructure} and
while \citet{Wang2022OnLearning} establish convergence to $\qopt$ when the optimal policy of the MDP is feedforward.
Overall, \citet[p. 99]{Sutton2018ReinforcementIntroduction} conjecture that under first-visit updates, sample averaging learning rates and exploring starts ``convergence to $\qopt$ seems inevitable''.
The following section discusses these results in detail.


%%%%%%%%%%%%%%%%%%%%%%%%%%%%%%%%%%%%%%%%%
%%%%%%%%%%%%%%%%%%%%%%%%%%%%%%%%%%%%%%%%%
\subsection{Counterexample with unconditional update distributions}
\label{sec: existing counter}

Example~5.12 in \citep{Bertsekas1996Neuro-DynamicProgramming} shows that under the modified Robbins-Monro conditions alone, solutions of the difference inclusion \eqref{eq: complete difference inclusion} can converge to a suboptimal fixed point on the boundary between the regions of several policies.

\begin{remark}[All suboptimal fixed points lie on boundaries]
    The sequence $(\pol_k)$ generated by MC-O-PI cannot converge to a suboptimal policy. If it did, the estimate $q_k$ would converge to the corresponding action values, which in turn would cause the policy to change. Therefore, any suboptimal fixed point that attracts the sequence $(q_k)$ must lie on the boundary between the regions of several policies so that the sequence $(\pol_k)$ can oscillate indefinitely.
\end{remark}

The counterexample presented below follows the strategy of Example~5.12 but differs in implementation.
We track an action-value function $q_k$ to compute the greedy policy at each step in a model-free manner,
whereas Bertsekas and Tsitsiklis track a state-value function $v_k$, which requires knowing the reward function $r$, the discount factor $\discount$ and the transition probabilities $\transit$ to compute the greedy policy.
Regardless, the two examples are equivalent if we define
\begin{align*}
    q_k(s,a) = r(s,a) + \discount \sum_{s_+ \in \sset} \transit(s_+ \mid s,a) v_k(s_+).
\end{align*}
This ensures that the greedy policies computed on $q_k$ and $v_k$ match.
% A similar construction appears in Appendix~F of \citep{Wang2022OnLearning}.


Concretely, consider the 
% discounted, single-state, deterministic 
MDP in \autoref{fig: experiment smallest MDP} with discount factor $\discount$. The state $s$ offers two actions: \textit{left} yields a reward of one and \textit{right} yields no reward. The best policy $\best{\pol}$ follows the \textit{left} action, denoted by $\best{a}$, while the worst policy $\worst{\pol}$ follows the \textit{right} action, denoted by $\worst{a}$. 
% For clarity, we denote the best action by $\best{a}$ and the suboptimal action by $\worst{a}$.
The action values of policies $\best{\pol}$ and $\worst{\pol}$ are:
\begin{align*}
    q_{\best{\pol}}(s, \best{a}) &= 1 + \discount + \discount^2 + \cdots = \dfrac{1}{1-\discount}
    \\
    q_{\best{\pol}}(s, \worst{a}) &= 0 + \discount + \discount^2 + \cdots = \dfrac{\discount}{1-\discount}
    \\
    q_{\worst{\pol}}(s, \best{a}) &= 1 + 0 + 0 + \cdots = 1
    \\
    q_{\worst{\pol}}(s, \worst{a}) &= 0 + 0 + 0 + \cdots = 0
\end{align*}

\begin{figure}[b]
    \centering
    \hrule
    \vspace{1.3cm}
    \begin{subfigure}[b]{0.3\textwidth}
        \centering
        \begin{tikzpicture}[
            scale=1, auto, node distance=2cm, 
            every edge/.style={draw, <-, font=\small},
            state/.style={circle, fill=mylightgrey, draw=black, text=black}
        ]
            % Nodes
            \node[state] (s1) at (0,0) {$s$};
            % Edges
            \path[->]
                (s1) edge[out=160, in=200, loop, looseness=8, draw=black] node[pos=0.5, left, text=black] {1} (s1)
                (s1) edge[out=20, in=340, loop, looseness=8, draw=black] node[pos=0.5, right, text=black] {0} (s1) ;
                % (s1) edge[out=145, in=215, loop, looseness=7] 
                %     node[pos=0.5, right] {$\best{a}$} 
                %     node[pos=0.5, left] {1} 
                % (s1)
                % (s1) edge[out=35, in=325, loop, looseness=7] 
                %     node[pos=0.5, left] {$\worst{a}$} 
                %     node[pos=0.5, right] {0} 
                % (s1);
        \end{tikzpicture}
        \caption{MDP}
        \label{fig: experiment smallest MDP}
    \end{subfigure}%
    \hfill
    \begin{subfigure}[b]{0.3\textwidth}
        \centering
        \begin{tikzpicture}[
            scale=1, auto, node distance=2cm, 
            every edge/.style={draw, <-, font=\small},
            state/.style={circle, fill=mylightgrey, draw=black, text=black}
        ]
            % Nodes
            \node[state] (s1) at (0,0) {$s$};
            % Edges
            \path[->]
                (s1) edge[out=160, in=200, loop, looseness=8, draw=black] node[pos=0.5, left, text=black] {1} (s1)
                (s1) edge[out=20, in=340, loop, looseness=8, draw=mylightgrey] node[pos=0.5, right, text=mylightgrey] {0} (s1) ;
                % (s1) edge[out=160, in=200, loop, looseness=8] node[pos=0.5, left] {$\begin{aligned}
                %     & \best{a} \\
                %     & 1
                % \end{aligned}$} (s1)
                % (s1) edge[out=20, in=340, loop, looseness=8, draw=mylightgrey] node[pos=0.5, right, text=mylightgrey] {$\begin{aligned}
                %     & \worst{a} \\
                %     & 0
                % \end{aligned}$} (s1);
        \end{tikzpicture}
        \caption{Policy $\best{\pol}$}
    \end{subfigure}%
    \hfill
    \begin{subfigure}[b]{0.3\textwidth}
        \centering
        \begin{tikzpicture}[
            scale=1, auto, node distance=2cm, 
            every edge/.style={draw, <-, font=\small},
            state/.style={circle, fill=mylightgrey, draw=black, text=black}
        ]
            % Nodes
            \node[state] (s1) at (0,0) {$s$};
            % Edges
            \path[->]
                (s1) edge[out=160, in=200, loop, looseness=8, draw=mylightgrey] node[pos=0.5, left, text=mylightgrey] {1} (s1)
                (s1) edge[out=20, in=340, loop, looseness=8, draw=black] node[pos=0.5, right, text=black] {0} (s1) ;
                % (s1) edge[out=160, in=200, loop, looseness=8, draw=mylightgrey] node[pos=0.5, left, text=mylightgrey] {$\begin{aligned}
                %     & \best{a} \\
                %     & 1
                % \end{aligned}$} (s1)
                % (s1) edge[out=20, in=340, loop, looseness=8] node[pos=0.5, right] {$\begin{aligned}
                %     & \worst{a} \\
                %     & 0
                % \end{aligned}$} (s1);
        \end{tikzpicture}
        \caption{Policy $\worst{\pol}$}
    \end{subfigure}%
    \vspace{0.7cm}
    \caption{MDP on which we apply the strategy of Example~5.12 in \citep{Bertsekas1996Neuro-DynamicProgramming}.}
    \label{fig: experiment smallest counter}
\end{figure}



% As discussed in \autoref{sec: mcopi behaviour}, we assume that the effective learning rates satisfy \autoref{asp: stochastic conditions on effective learning rate}.

%%%%%%%%%%%%%%%%%%%%%%%%%%%%%%%%%%%%%%%%%
\subsubsection{Initial conditions}

\autoref{fig: smallest counter initial condition} depicts the action values $\qopt$ and $q_{\worst{\pol}}$ within the space $\qset$. The dotted line separates the regions $\graval(\best{\pol})$ and $\graval(\worst{\pol})$.
Thus points above the dotted line are attracted towards $q_{\worst{\pol}}$, while those below move towards $\qopt$.
As shown in \autoref{thm: bounded limit sets}, it suffices to consider solutions starting in $\qset_\infty$ when studying the asymptotic behaviour of MC-O-PI.

Since the MDP is deterministic, observed returns computed at the end of each episode match expected returns.
Thus the estimate $q_k$ moves closer to $q_{\pol_k}$ at every step for at least one state-action pair without overshooting it.
As a result, every solution that starts in an initial condition $q_0 \in \qset_\infty$ satisfying $q_0(s, \best{a}) \geq q_{\best{\pol}}(s, \worst{a})$, always remains in the region $\region({\best{\pol}})$ and converges to $\qopt$.
Conversely, solutions that do not converge to $\qopt$ start in the region
\begin{align}
    \qset_0 \coloneqq \big\{ q \in \qset_\infty : q(s, \best{a}) < q_{\best{\pol}}(s, \worst{a}) \big\} .
\end{align}
As illustrated in \autoref{fig: smallest counter initial condition}, $\qset_0$
% overlaps the regions $\region_{\best{\pol}}$ and $\region_{\worst{\pol}}$.
is not empty provided
\begin{align*}
    q_{\worst{\pol}}(s, \best{a}) = 1 
    < \dfrac{\discount}{1-\discount} = q_{\best{\pol}}(s, \worst{a}) ,
\end{align*}
which is satisfied as long as $\discount > 1/2$.

\begin{figure}[hb!]
    \centering
    \hrule
    \vspace{0.4cm}
    \begin{subfigure}{0.48\textwidth}
        \centering
        \begin{tikzpicture}[scale=1.4]

                \filldraw[fill=CambridgeSuperLightBlue, draw=none] (1,0) rectangle (2,2);
                
                \draw[->] (0,0) -- (3.2,0) node[below, font=\small] {$q(s,\best{a})$};
                \draw[->] (0,0) -- (0,2.2) node[above, font=\small] {$q(s,\worst{a})$};

                \draw[dotted, black] (0,0) -- (2.2,2.2);
                
                \filldraw[fill=none, draw=black] (1,0) rectangle (2.5,2);
    
                % \draw[black] (2,2) -- (2,1.25);
                % \draw[black] (2,0) -- (2,0.75);
                
                \node at (1.5, 1) [font=\small, color=CambridgeCoreBlue] {$\region_0$};
                \node at (2,1) [font=\small, color=black] {$\region_\infty$};
        
                % Add labels to the axes
                \draw (2.5,0.1) -- (2.5,-0.1) node[below, font=\small] {$\frac{1}{1-\discount}$};
                \draw (0.1,2) -- (-0.1,2) node[left, font=\small] {$\frac{\discount}{1-\discount}$};
                \draw (1,0.1) -- (1,-0.1) node[below, font=\small] {$1$};
                \draw (0,0.1) -- (0,-0.1) node[left, font=\small, xshift=-4.5pt, yshift=-6.5pt] {0};
                \draw (-0.1,0) -- (0.1,0) node[font=\small] {};
    
                % Add the points with associated text
                \filldraw[fill=black, draw=black] (2.5,2) circle (1pt) node[above, text=black] {$q_{\best{\pol}}$};
                \filldraw[fill=black, draw=black] (1,0) circle (1pt) node[above left, text=black] {$q_{\worst{\pol}}$};
            
            \end{tikzpicture}
        \caption{Initial condition}
        \label{fig: smallest counter initial condition}
    \end{subfigure}
    \hfill
    \begin{subfigure}{0.48\textwidth}
        \centering
        \begin{tikzpicture}[scale=1.4]
                \draw[->] (0,0) -- (3.2,0) node[below, font=\small] {$q(s,\best{a})$};
                \draw[->] (0,0) -- (0,2.2) node[above, font=\small] {$q(s,\worst{a})$};

                \draw[dotted, black] (0,0) -- (2.2,2.2);
                
                \filldraw[fill=mylightgrey, draw=black] (1,0) rectangle (2,2);
                \draw[black] (1.4,0) -- (1.4,2);
                \draw[black] (1,1) -- (2,2);
                
                \node at (1, 0.5) [left, font=\small] {$\region_1$};
                \node at (2, 0.5) [right, font=\small] {$\region_2$};
                \node at (1.7, 2) [above, font=\small] {$\region_3$};
                \node at (1, 1.5) [left, font=\small] {$\region_4$};

                % Add labels to the axes
                \draw (2.5,0.1) -- (2.5,-0.1) node[below, font=\small] {$\frac{1}{1-\discount}$};
                \draw (0.1,2) -- (-0.1,2) node[left, font=\small] {$\frac{\discount}{1-\discount}$};
                \draw (1,0.1) -- (1,-0.1) node[below, font=\small] {$1$};
                \draw (0,0.1) -- (0,-0.1) node[left, font=\small, xshift=-4.5pt, yshift=-6.5pt] {0};
                \draw (-0.1,0) -- (0.1,0) node[font=\small] {};
    
                % Add the points with associated text
                \filldraw[fill=black, draw=black] (2.5,2) circle (1pt) node[above, text=black] {$q_{\best{\pol}}$};
                \filldraw[fill=black, draw=black] (1,0) circle (1pt) node[above left, text=black] {$q_{\worst{\pol}}$};
                \filldraw[fill=black, draw=black] (1.4, 1.4) circle (1pt) node[below right, text=black, xshift=-2.5pt] {$\hat{q}$};
                
                \draw[->, CambridgeCoreBlue, thick] (1.2, 0.4) -- (1.8, 0.4) -- (1.8, 1.9) -- (1.2, 1.9) -- (1.2, 0.8) -- (1.6, 0.8) -- (1.6, 1.7) -- (1.45, 1.7) ;
                
            \end{tikzpicture}
        \caption{Updating strategy}
        \label{fig: smallest counter update}
    \end{subfigure}
    \caption{Solutions of MC-O-PI can converge to any point $\hat{q}$ along the boundary in $\qset_0$ when starting in $\qset_0$ and using the updating strategy described on page \pageref{enum: existing update distrib}.}
    \label{fig: experiment for unconditional bias}
\end{figure}



%%%%%%%%%%%%%%%%%%%%%%%%%%%%%%%%%%%%%%%%%
\subsubsection{Updating strategy}
\label{sec: sampling strat of ex 512}


Take any point $\hat{q} \in \qset_0$ on the boundary between the regions $\region({\best{\pol}})$ and $\region({\worst{\pol}})$.
% \begin{align*}
%     q_{\worst{\pol}}(s, \best{a}) < \hat{q}(s, \best{a}) = \hat{q}(s, \worst{a}) < q_{\best{\pol}}(s, \worst{a})
% \end{align*}
Further, let $\qset_1$, $\qset_2$, $\qset_3$ and $\qset_4$ denote the four regions that result from splitting $\qset_0$ along the boundary and vertically through $\hat{q}$.
We then obtain the spiral behaviour depicted in \autoref{fig: smallest counter update}, by combining initial-visit updates with the following sampling strategy:
\begin{enumerate}

    \item \label{enum: existing update distrib}
    When $q_\kpz \in \qset_1$, start each episode in $(s, \best{a})$.
    The solution then moves horizontally towards $q_{\best{\pol}}$ until it enters $\qset_2$.
    
    \item 
    When $q_\kpz \in \qset_2$, start each episode in $(s, \worst{a})$. 
    The solution then moves vertically towards $q_{\best{\pol}}$ until it crosses the boundary between the regions $\region({\best{\pol}})$ and $\region({\worst{\pol}})$,
    and enters $\qset_3$.

    \item
    When $q_\kpz \in \qset_3$, start each episode in $(s, \best{a})$. 
    The solution then moves horizontally towards $q_{\worst{\pol}}$ until it enters $\qset_4$.

    \item
    When $q_\kpz \in \qset_4$, start each episode in $(s, \worst{a})$. 
    The solution then moves vertically towards $q_{\worst{\pol}}$ until it crosses the boundary between the regions $\region({\worst{\pol}})$ and $\region({\best{\pol}})$,
    and enters $\qset_1$.

\end{enumerate}
% definition of \updateste
% Thus, at every step the update set $\updateset_k$ is defined over $\qset_0$ as 
% \begin{align}
%     \updateset_k(q) =
%     \begin{cases}
%         \ \update_k(s, \best{a})=1, \ \update_k(s, \worst{a})=0 &q \in \qset_1
%         \\
%         \ \update_k(s, \best{a})=0, \ \update_k(s, \worst{a})=1 &q \in \qset_2
%         \\
%         \ \update_k(s, \best{a})=1, \ \update_k(s, \worst{a})=0 &q \in \qset_3
%         \\
%         \ \update_k(s, \best{a})=0, \ \update_k(s, \worst{a})=1 &q \in \qset_4
%     \end{cases}
% \end{align}

% bound stepsize
As the sequence of learning rates decreases, the estimate $q_k$ converges to $\hat{q}$ while the policy $\pol_k$ oscillates indefinitely between $\best{\pol}$ and $\worst{\pol}$.
The learning rates should be sufficiently small to ensure that solutions in $\qset_1$ do not overshoot $\qset_2$ and leave $\qset_0$. Thus when $q_\kpz \in \qset_1$ we require that
\begin{align*}
    q_\kpo(s, \best{a}) = q_k(s, \best{a}) + \lrate_k (q_{\best{\pol}}(s, \best{a}) - q_k(s, \best{a})) < q_{\best{\pol}}(s, \worst{a}) ,
\end{align*}
which is satisfied as long as
\begin{align*}
    \lrate_k
    % < \dfrac{q_{\best{\pol}}(s, \worst{a}) - q_k(s, \best{a})}{q_{\best{\pol}}(s, \best{a}) - q_k(s, \best{a})}
    < \dfrac{q_{\best{\pol}}(s, \worst{a}) - \hat{q}(s, \worst{a})}{q_{\best{\pol}}(s, \best{a}) - \hat{q}(s, \best{a})}
\end{align*}
for every $k \geq 0$.


%%%%%%%%%%%%%%%%%%%%%%%%%%%%%%%%%%%%%%%%%
\subsubsection{Discussion}


This counterexample shows that under the modified Robbins-Monro conditions alone, MC-O-PI may converge to a suboptimal fixed point, even for an MDP with only two deterministic policies.

This example is unrealistic for several reasons. First, it assumes we know $\qopt$ and use it to carefully select initial state-action pairs and learning rates that keep solutions in $\qset_0$.
% In practice, $\qopt(s, \worst{a})$ is unknown, making it impossible to carefully select initial state-action pairs to prevent solutions from leaving $\qset_0$.
If instead we sample initial state-action pairs from a distribution that satisfies exploring starts, there is a non-zero probability that solutions leave $\qset_0$. For instance, if $q_k \in \qset_2$, several updates of $q_k(s, \best{a})$ would drive solutions out of $\qset_0$.

Furthermore, this example relies on exact policy evaluation to avoid overshooting $\qset_2$. However, in an MDP with stochastic transitions, any overestimation of $\qopt$ could push solutions out of $\qset_0$.

Finally, this example applies initial-visit updates which are less data-efficient than first-visit updates; an important aspect in practical applications. 
As soon as we implement first-visit updates, the spiral breaks because solutions in $\qset_2$ move along a straight line towards $\qopt$ and thus never leave $\region(\best{\pol})$.

These observations raise the question: Does a counterexample exist where
initial state-action pairs satisfy exploring starts, 
episode simulations are stochastic 
and the algorithm applies first-visit updates?
Under these conditions, each update distribution $\update_k$ is independent of $q_k$'s position in $\qset$ but may depend on the currently greedy policy $\pol_k$.
% as when using first-visit updates or an epsilon-greedy sampling strategy for initial state-action pairs.

% allow update distribution to change when policy changes as this is often the case
% - eg epsilon-greedy sampling
% - first visit updates

% %%%%%%%%%%%%%%%%%%%%%%%%%%%%%%%%%%%%%%%%%
% \paragraph{Example breaks down with stochastic transitions}

% % stochastic transition can only come from episodic
% The deterministic MDP described in \autoref{fig: experiment smallest counter} allows for exact policy evaluation. However, it requires infinitely long episodes, which is infeasible in practice.
% To address this problem, let us terminate the episode after each action with probability $\delta$.
% The action-values of policies $\best{\pol}$ and $\worst{\pol}$ then become:
% \begin{align*}
%     &q_{\best{\pol}}(s, \best{a}) = 1 + \discount(1-\delta) + \discount^2(1-\delta)^2 + \cdots = \dfrac{1}{1-\discount(1-\delta)}
%     \\
%     &q_{\best{\pol}}(s, \worst{a}) = 0 + \discount(1-\delta) + \discount^2(1-\delta)^2 + \cdots = \dfrac{\discount(1-\delta)}{1-\discount(1-\delta)}
%     \\
%     &q_{\worst{\pol}}(s, \best{a}) = 1 + 0 + 0 + \cdots = 1
%     \\
%     &q_{\worst{\pol}}(s, \worst{a}) = 0 + 0 + 0 + \cdots = 0
% \end{align*}
% and the region $\qset_0$ is not empty as long as $\discount(1-\delta) > 1/2$.


% % stochastic transitions mean noisy policy evaluation
% % The stochastic transitions introduce noise in the observed returns.
% % to be noisy estimates of the action values.
% The length of an episode is now $1/\delta$, on average.
% However, due to the stochastic transitions, any episode can be much longer than expected, leading to an overestimation of the action values and a large update of $q_k$.
% As a result, it becomes impossible to keep a solution of MC-O-PI inside $\qset_0$.
% For instance, when $q_k \in \qset_1$ the solution overshoots $\qset_2$ and leaves $\qset_0$ if
% \begin{align*}
%     q_\kpo(s, \best{a})
%     = q_k(s, \best{a}) + \lrate_k \big( q_{\best{\pol}}(s, \best{a}) - q_k(s,\best{a}) + w_k(s, \best{a}) \big) \geq q_{\best{\pol}}(s, \worst{a}) ,
% \end{align*}
% or, in other words, if
% \begin{align*}
%     w_k(s, \best{a})
%     \geq
%     \dfrac{1}{\lrate_k} \big( q_{\best{\pol}}(s, \worst{a}) - q_k(s,\best{a}) \big) - \big( q_{\best{\pol}}(s, \best{a}) - q_k(s,\best{a}) \big) .
% \end{align*}
% Thus even if the learning rate $\lrate_k$ is small and the solution is close to $\hat{q}$, the update could overshoot $\qset_2$, leave $\qset_0$ and converge to $\qopt$.
% Although the factor $1/\lrate_k$ indicates that the probability of this happening decreases to zero as $k$ tends to infinity, given that the noise has a zero mean and a bounded variance.

% computing the probability of this happening once
% if is finite => prob of staying is non-zero


% %%%%%%%%%%%%%%%%%%%%%%%%%%%%%%%%%%%%%%%%%
% \paragraph{Example breaks down with first-visit updates}

% If we replace initial-visit updates with first-visit updates, then the update distribution $\update$ becomes uniform when the estimate $q_k$ lies in region $\qset_2$.
% % or $\qset_3$.
% As a result, from any point in $\qset_2$, the solution moves, on average, along a straight line towards $\qopt$. 
% Since no straight line leads to region $\qset_3$, the spiral behaviour is not sustainable on the long term. Instead, all solutions converge to $\qopt$.

% % Wang and Ross example is not applicable here bc they visit best region (we show that with greedy bias when enter best region leave the boundary)
% % - need to repeatedly sample the off-policy action


%%%%%%%%%%%%%%%%%%%%%%%%%%%%%%%%%%%%%%%%%
%%%%%%%%%%%%%%%%%%%%%%%%%%%%%%%%%%%%%%%%%
\subsection{Conjecture of Sutton and Barto}

The previous example highlights that additional statistical
regularity beyond the Robbins-Monro conditions is necessary to guarantee that solutions of MC-O-PI converge to $\qopt$.
\citet[p. 99]{Sutton2018ReinforcementIntroduction}
conjecture that first-visit updates, sample averaging learning rates and exploring starts
suffice to ensure convergence.
However, no proof or counterexample has settled this claim so far.

%%%%%%%%%%%%%%%%%%%%%%%%%%%%%%%%%%%%%%%%%
%%%%%%%%%%%%%%%%%%%%%%%%%%%%%%%%%%%%%%%%%
\subsection{Convergence proof for uniform update distributions}

\vspace{-0.1cm}

Suppose all state-action pairs are updated equally frequently on average. The update distributions are then uniform over all state-action pairs.

\vspace{-0.1cm}

\begin{definition}[Uniform updates over all state-action pairs]
\label{asp: unbiased tsitsiklis}
The update distribution $\update$ is uniform over all state-action pairs
if and only if
for all states $s, s' \in \sset$ and actions $a \in \aset(s)$ and $a' \in \aset(s')$, we have
\begin{align}
    \nonumber
    \update(s, a) = \update(s', a') .
\end{align}
\end{definition}

\vspace{-0.7cm}

\cite{Tsitsiklis2002OnIteration} establishes that MC-O-PI converges to the optimal solution under this uniformity assumption if the discount factor $\discount<1$.

% similar to the main site of the proof of modified policy iteration \cite{Puterman1978ModifiedProblems}

\vspace{-0.1cm}

\begin{theorem}
\label{thm: tsitsiklis convergence proof}
Solutions of MC-O-PI (\autoref{def: mcopi}) converge almost surely to $q_{\best{\pol}}$
if the update distributions $(\update_k)_{k \geq 0}$ are uniform over all state-action pairs (\autoref{asp: unbiased tsitsiklis}),
the learning rates $(\lrate_k)_{k \geq 0}$ and update distributions $(\update_k)_{k \geq 0}$ satisfy the modified Robbins-Monro conditions (\autoref{asp: stochastic conditions on effective learning rate})
and the discount factor $\discount < 1$.
\end{theorem}

\vspace{-0.7cm}

The proof of \autoref{thm: tsitsiklis convergence proof} relies on the following operators:

\vspace{-0.1cm}

\begin{definition}[Bellman operators]
\label{def: bellman operators}
The transition operators $\T_\pol, \T : \qset \to \qset$ are defined as
\begin{align*}
    (\T_{\pol} q)(s,a) &\coloneqq \sum_{s_+ \in \sset} \transit(s_+ \mid s,a) q(s_+,\pol) ,
    \\
    (\T q)(s,a) &\coloneqq \sum_{s_+ \in \sset} \transit(s_+ \mid s,a) \max_{a_+ \in \aset(s_+)} q(s_+, a_+) .
\end{align*}
The Bellman operators $\B_\pol, \B : \qset \to \qset$ are defined as
\begin{align*}
    (\B_\pol q) (s,a) &\coloneqq \reward(s,a) + \discount \sum_{s_+ \in \sset} \transit(s_+ \mid s,a) q(s_+, \pol) ,
    \\
    (\B q)(s,a) &\coloneqq \reward(s,a) + \discount \sum_{s_+ \in \sset} \transit(s_+ \mid s,a) \max_{a_+ \in \aset(s_+)} q(s_+,a_+) .
\end{align*}
\end{definition}

When $\discount<1$, the Bellman operators are contracting and monotone \citep[Example 1.2.1]{Bertsekas2018AbstractProgramming}. These properties imply that the action values $q_\pol$ 
% and $\qopt$ 
are the unique solutions of the linear system
\begin{align}
\label{eq: linear system with operators}
    q_\pol &= r + \discount \T_{\pol} q_\pol = \B_{\pol} q_\pol
    % \\
    % \qopt &= r + \discount \T \qopt = \B \qopt
\end{align}
and also that, for every $q \in \qset$, if $q \leq \B_{\pol} q$
% or $q \leq \B q$, 
it follows that
\begin{align*}
    q \leq \B_{\pol} q \leq \B_{\pol}^2 q \leq \dots \leq \lim_{n \to \infty} \B_{\pol}^n q = q_\pol .
    % \\
    % q \leq \B q \leq \B^2 q \leq \dots \leq \lim_{n \to \infty} \B^n q = \qopt .
\end{align*}
% \begin{property}[Monotonicity and contraction of Bellman operators] 
% \label{thm: bellman operator monotonicity}
Every Bellman operator is monotone, meaning that for every policy $\pol$ and all $q, q' \in \qset$ it holds that
\begin{align*}
    q \leq q' \implies \B_\pol q \leq \B_\pol q' .
\end{align*}
% \end{property}
% 
% \begin{property}[Contraction in discounted MDP] 
% \label{thm: bellman operator contraction}
Every Bellman operator in a discounted MDP is also contracting, meaning that 
% if the discount factor $\discount<1$, then 
for every policy $\pol$ and all $q, q' \in \qset$ it holds that
\begin{align*}
    \|q - q'\| \geq \discount \|\B_\pol q - \B_\pol q'\| .
\end{align*}
\end{property}
% 
% Under the conditions of \autoref{thm: tsitsiklis convergence proof}, solutions of \eqref{eq: complete difference inclusion} satisfy
The key insight for proving \autoref{thm: tsitsiklis convergence proof} is that, with uniform update distributions over all state-action pairs, solutions of \eqref{eq: complete difference inclusion} satisfy
\begin{align}
\label{eq: tsitsiklis central result}
    \limsup_{k \to \infty} q_k - \B q_k \leq \zerofun .
\end{align}
As a result, for sufficiently large $k$, we obtain 
\begin{align*}
    q_k \leq \B q_k = \B_{\pol_k} q_k \leq \B_{\pol_k}^2 q_k \leq \dots \leq q_{\pol_k} .
\end{align*}
Thus solutions converge because they eventually tend to increase monotonically and are bounded above.

When $\discount = 1$, the Bellman operators may not be contracting, but \cite{Chen2018OnProblem} and \cite{Liu2021OnStarts} extend Tsitsiklis' framework under additional assumptions that ensure bounded optimal action values.
% Since their arguments are not essential to our analysis, we do not further discuss their proofs.
% 
To emphasise the importance of the uniformity assumption in the proofs of Tsitsiklis, Chen and Liu, we derive expression \eqref{eq: tsitsiklis central result}, but refer to \citep[p.~63-66]{Tsitsiklis2002OnIteration} for the remaining technical details.
% of \autoref{thm: tsitsiklis convergence proof}.

% \begin{lemma}
\label{thm: tsitsiklis helper lemma}
For every deterministic policy $\pol \in \polset$
and every point $q \in \region(\pol)$ it holds that
\begin{align}
\label{eq: property of c(q)}
    q_\pol \geq \B q - \dfrac{\discount c(q)}{1-\discount} \onefun .
\end{align}
\end{lemma}
\begin{proof}
By definition, $\B q = \B_\pol q$. Thus, equation \eqref{eq: def c tsitsiklis} implies that
\begin{align*}
    \B_\pol q \geq q - c(q) \onefun .
\end{align*}
Applying the monotone operator $\B_\pol$ to both sides of the inequality yields
\begin{align*}
    \B_\pol^2 q 
    &\geq \B_\pol (q - c(q) \onefun)
    = \B_\pol q - \discount c(q) \onefun
    \geq q - c(q) \onefun - \discount c(q) \onefun .
\end{align*}
Continuing inductively, we obtain that
\begin{align*}
    \B_\pol^n q \geq q - (1 + \discount + \dots + \discount^{n-1}) c(q) \onefun
\end{align*}
for all $n \geq 1$.
As $n$ tends to infinity, $\B_\pol^n q$ converges to $q_\pol$ and $(1 + \dots + \discount^{n-1})$ converges to $1/(1-\discount)$. Since the limit is well-defined for both sides of the inequality, we observe that
\begin{align*}
    q_\pol \geq q - \dfrac{c(q)}{1-\discount} \onefun .
\end{align*}
Finally, applying the operator $\B_\pol$ once more to both sides of this inequality yields equation \eqref{eq: property of c(q)}.
\end{proof}
% The second part of Tsitsiklis' proof contains the key insight.

\begin{lemma}
\label{thm: c(q) goes to zero}
If the update distributions $(\update_k)_{k \geq 0}$ are uniform over all state-action pairs (\autoref{asp: unbiased tsitsiklis}),
and together with the learning rates $(\lrate_k)_{k \geq 0}$ they satisfy the modified Robbins-Monro conditions (\autoref{asp: stochastic conditions on effective learning rate}),
an MC-O-PI solution $(q_k)_{k \geq 0}$ (\autoref{def: mcopi}) satisfies almost surely
\begin{align}
\label{eq: condition tsitsiklis eventual contract}
    \limsup_{k \to \infty} \ q_k - \B q_k \leq \zerofun .
\end{align}
\end{lemma}
\begin{proof}
    Let $(q_k)_{k \ge 0}$ and $(\pol_k)_{k \ge 0}$ be solutions of \eqref{eq: complete difference inclusion} with 
    % initial condition $q_0$, 
    learning rates $(\lrate_k)_{k \ge 0}$, update distributions $(\update_k)_{k \ge 0}$ and perturbations $(w_k)_{k \ge 0}$.
    The affine properties of the Bellman operators yield
    \begin{align}
        \nonumber
        \B q_\kpo &= \B_{\pol_\kpo} q_\kpo
        \\
        \nonumber
        &\geq \B_{\pol_k} q_\kpo
        \\
        \nonumber
        &= \B_{\pol_k} \big( (\I- \lrate_k \update_k) q_k + \lrate_k \update_k (q_{\pol_k}+w_k) \big)
        \\
        \label{eq: tsitsiklis commutation}
        &= r + (\I-\lrate_k \update_k) \discount \T_{\pol_k} q_k + \lrate_k \update_k \discount \T_{\pol_k} ( q_{\pol_k} + w_k)
        \\
        \nonumber
        &= (\I-\lrate_k \update_k) \B_{\pol_k} q_k + \lrate_k \update_k \B_{\pol_k} q_{\pol_k} + \lrate_k \discount \T_{\pol_k} \update_k w_k
        \\
        \nonumber
        &= (\I-\lrate_k \update_k) \B q_k + \lrate_k \update_k q_{\pol_k} + \lrate_k \discount \T_{\pol_k} \update_k w_k
        % \\
        % \nonumber
        % &= (\I-\lrate_k \update_k) (\B q_k - q_k + q_k) + \lrate_k \update_k ( q_{\pol_k} + w_k - w_k) + \lrate_k \discount \T_{\pol_k} \update_k w_k
        \\
        \nonumber
        &= (\I-\lrate_k \update_k) (\B q_k - q_k) + q_\kpo + \lrate_k (\discount \T_{\pol_k} \update_k w_k -  \update_k w_k) 
    \end{align}
    As a result, we obtain
    \begin{align}
    \label{eq: q-Bq decreases}
        q_\kpo - \B q_\kpo \leq (\I - \lrate_k \update_k)(q_k - \B q_k) - \lrate_k (\discount \T_{\pol_k} \update_k w_k -  \update_k w_k) .
    \end{align}
    % where
    % \begin{align*}
    %     &\E \big[ \discount \T_{\pol_k} \update_k w_k - \update_k w_k \big] = \zerofun 
    %     \\
    %     &\E \big[ \| \discount \T_{\pol_k} \update_k w_k - \update_k w_k \|^2 \big] \leq c_1 + c_2 \| q_k \|^2
    % \end{align*}
    % for some constants $c_1, c_2 \in \real_{\geq 0}$ independent of $k$.
    Sections 4.2 and 4.3 in \citep{Bertsekas1996Neuro-DynamicProgramming} show that since the sequence $(q_k - \B q_k)_{k \geq 0}$ satisfies \eqref{eq: q-Bq decreases}, it is upper-bounded by a sequence that converges to zero with probability one. As a result, we obtain almost surely expression \eqref{eq: condition tsitsiklis eventual contract}.
    % \begin{align*}
    %     \limsup_{k \to \infty} q_k - \B q_k \leq 0 .
    % \end{align*}
\end{proof}

% We can now prove Tsitsiklis' main result.
% % \begin{theorem}
% \label{thm: tsitsiklis convergence proof}
% If (1) the effective learning rates $(\lrate_k \update_k)_{k \geq 0}$ satisfy the Robbins-Monro conditions (\autoref{asp: stochastic conditions on effective learning rate}),
% % If the learning rates $(\lrate_k)$ and the update maps $(\updateset_k)$ satisfy \autoref{asp: stochastic conditions on effective learning rate},
% (2) all update distributions $(\update_k)_{k \geq 0}$ are uniform (\autoref{asp: unbiased tsitsiklis}),
% and (3) the discount factor satisfies $0 \leq \discount < 1$,
% then every sequence $(q_k)_{k \geq 0}$ generated by MC-O-PI converges to $q_{\best{\pol}}$ with probability one.
% \end{theorem}

\begin{proof}[Proof of \autoref{thm: tsitsiklis convergence proof}]
Consider any sequence $(q_k)$ generated by MC-O-PI and let $(\pol_k)$, $(\lrate_k \update_k)$ and $(w_k)$ be the associated policies, effective learning rates and noise terms.
\autoref{thm: c(q) goes to zero} indicates that for every $\epsilon > 0$ there exists a $K$ such that $c(q_k) \leq \epsilon$ for all $k \geq K$. Consequently, \autoref{thm: tsitsiklis helper lemma} implies that
\begin{align*}
    q_\kpo
    &= (\I - \lrate_k \update_k) q_k + \lrate_k \update_k (q_{\pol_k} + w_k)
    \\
    & \geq (\I - \lrate_k \update_k) q_k + \lrate_k \update_k \Big( \B q_k - \frac{\discount \epsilon}{1 - \discount} \onefun + w_k \Big) .
\end{align*}
for all $k \geq K$.
Thus, consider the system
\begin{align*}
    x_\kpo = (\I - \lrate_k \update_k) q_k + \lrate_k \update_k \Big( \B q_k - \frac{\discount \epsilon}{1 - \discount} \onefun + w_k \Big)
\end{align*}
% for $k \geq K$ 
initialised as $x_K = q_K$. Clearly, $q_k \geq x_k$ for all $k \geq K$.
Proposition 4.4 in \citep{Bertsekas1996Neuro-DynamicProgramming} then indicates that the sequence $(x_k)_{k \geq K}$ converges with probability one to
\begin{align*}
    q_{\best{\pol}} - \frac{\discount \epsilon}{(1-\discount)^2} \onefun .
\end{align*}
Since $\epsilon$ can be chosen arbitrarily close to zero, we conclude that
\begin{align*}
    \liminf_{k \to \infty} q_k \geq \sup_{\epsilon > 0}
    q_{\best{\pol}} - \frac{\discount \epsilon}{(1-\discount)^2} \onefun
    =
    q_{\best{\pol}} .
\end{align*}
Finally, \autoref{thm: bounded solution} indicates that
\begin{align*}
    \limsup_{k \to \infty} q_k \leq q_{\best{\pol}}
\end{align*}
which completes the proof.
\end{proof}

\vspace{-0.3cm}

The proofs of \autoref{thm: c(q) goes to zero} and \autoref{thm: tsitsiklis convergence proof} rely on the operator $\T_{\pol_k}$ commuting with the effective learning rate $\lrate_k \update_k$ in \eqref{eq: tsitsiklis commutation}.
This only holds when $\lrate_k \update_k$ is uniform over all state-action pairs.

\begin{remark}[Only the effective learning rates in the limit matter]
\label{rmk: eff lr in the limit}
    Tsitsiklis notes that MC-O-PI's asymptotic behaviour depends on the asymptotic behaviour of the effective learning rates $(\lrate_k \update_k)$, not on either component alone.
    For instance, if the update distributions are non-uniform but constant, say $\update_k = \update$ for all $k$, and the learning rates take the form $\lrate_k(s,a) = 1/n_k(s,a)$, where $n_k(s,a)$ denotes the number of updates of $(s,a)$ in the first $k$ steps,
    then $\lrate_k \update_k$ converges to $1/k$ because $n_k(s,a)$ tends to $k \update(s,a)$ for each state-action pair. As a result, \autoref{thm: c(q) goes to zero} still applies.
    % This highlights that MC-O-PI's asymptotic behaviour depends on the asymptotic behaviour of the effective learning rates $(\lrate_k \update_k)$, rather than either component individually.
\end{remark}


\begin{remark}[Lyapunov interpretation]
Let $(q)_+ (s,a) = \max \{ 0, q(s,a) \}$ denote the componentwise ReLU operator,
and consider the potential function that measures the one-sided Bellman residual
\begin{align}
    V(q) = \frac{1}{2} \|(q-\B q)_+\|_2^2 .
\end{align}
The optimal action-value function $\qopt$ is clearly an equilibrium of this potential.
Using the linear system defined in \eqref{eq: linear system with operators}, we obtain the gradient
\begin{align}
    \nabla V (q) = (\I - \discount \T_\pi )^\top (q-\B q)_+
\end{align}
where $\pi$ is any deterministic policy that is greedy with respect to $q$.
Computing the inner product with the mean-field update of MC-O-PI when updates are uniform over all state-action pairs yields
\begin{align*}
    \langle \nabla V(q) , q_\pi - q \rangle
    &= \langle (\I - \discount \T_\pi )^\top (q-\B q)_+ , q_\pi - q \rangle
    \\
    &= \langle (\I - \discount \T_\pi )^\top (q-\B q)_+ , (\I - \discount \T_\pi )^{-1} (\B q - q) \rangle
    \\
    &= \langle (q-\B q)_+ , (\B q - q) \rangle
    \\
    &= - \|(q - \B q)_+\|_2^2
    \\
    &= - 2 V(q)
\end{align*}
As a result, when updates are uniform over all state-action pairs, the mean-field MC-O-PI update is a descent direction for the potential $V$.
This geometric insight corresponds to the decrease in inequality \eqref{eq: q-Bq decreases} that is essential to Tsitsiklis' proof.

Overall, this Lyapunov interpretation offers an alternative method to prove \autoref{thm: tsitsiklis convergence proof}, using Proposition~4.1 of \citep{Bertsekas1996Neuro-DynamicProgramming}.
\end{remark}


%%%%%%%%%%%%%%%%%%%%%%%%%%%%%%%%%%%%%%%%%
%%%%%%%%%%%%%%%%%%%%%%%%%%%%%%%%%%%%%%%%%
\subsection{Convergence proof for feedforward optimal policy}

\citet{Wang2022OnLearning} show that MC-O-PI converges to the optimal solution if no state is visited twice when simulating an episode with policy $\best{\pol}$.

% The optimal policy is then feedforward.

\vspace{-0.2cm}

\begin{definition}[Feedforward optimal policy]
\label{asp: wang feedforward}
    The optimal policy $\best{\pol}$ is feedforward if the Markov chain induced by this policy on the MDP is acyclic.
\end{definition}

\vspace{-0.3cm}

\begin{theorem}
\label{thm: optimal feedforward}
Solutions of MC-O-PI (\autoref{def: mcopi}) converge almost surely to $\qopt$
if 
the learning rates $(\lrate_k)_{k \geq 0}$ and update distributions $(\update_k)_{k \geq 0}$ satisfy the modified Robbins-Monro conditions (\autoref{asp: stochastic conditions on effective learning rate})
and the optimal policy is feedforward (\autoref{asp: wang feedforward}).
\end{theorem}

\vspace{-0.5cm}

The key insight for proving \autoref{thm: optimal feedforward} is that the feedforward structure of $\best{\pol}$ induces an order on the state space where each state-action pair $(s_t, \best{\pol})$ transitions only to states $s_{t'}$ with $t' > t$.
As a result, if $\pol_k$ has converged to $\best{\pol}$ in all states $(s_{t'})_{t' > t}$, every episode that visits $(s_t, \best{\pol})$ yields an expected return of $\qopt(s_t, \best{\pol})$,
and consequently $q_k(s_t, \best{\pol})$ converges to $\qopt(s_t, \best{\pol})$.
Since $q_k(s_t, a)$ for suboptimal actions is bounded above by $\qopt(s_t, a) < \qopt(s_t, \best{\pol})$, $\pol_k$ is eventually greedy in $s_t$ as well.
Thus, with a finite state space of size $T$, MC-O-PI converges to the optimal solution, starting in $s_T$, then in $s_{T-1}$, and so on.
See \cite[Appendix~B]{Wang2022OnLearning} for the remaining technical details.

% \begin{proof}[Proof of \autoref{thm: optimal feedforward}]
The feedforward structure of the optimal policy $\best{\pol}$ induces an order on the finite state space.
In this order, each state-action pair $(s_n, \best{\pol})$ can only transition to states $s_{n'}$ where $n' > n$.

Consider the last state in the order, denoted by $s_N$.
By definition, every episode that visits $(s_N, \best{\pol})$ experiences a reward $r(s_N, \best{\pol})$ and terminates.
Since this reward corresponds to the action value $q_{\best{\pol}}(s_N, \best{\pol})$, the estimate is updated after every episode as
\begin{align*}
    q_\kpo(s_N, \best{\pol})
    =
    q_k(s_N, \best{\pol})
    + \lrate_k \update_k ( q_{\best{\pol}}(s_N, \best{\pol}) - q_k(s_N, \best{\pol}) + w_k(s_N, \best{\pol}) ) .
\end{align*}
The Robbins-Monro conditions guarantee that this update is performed infinitely often. As a result, $q_k(s_N, \best{\pol})$ converges to $q_{\best{\pol}}(s_N, \best{\pol})$ with probability one \citep[Example 4.3]{Bertsekas1996Neuro-DynamicProgramming}.
For all other actions, \autoref{thm: bounded solution} shows that the estimate $q_k(s_N, a)$ converges to a number that is at most $q_{\best{\pol}}(s_N, a)$. Since those actions are suboptimal, there exists an $\epsilon>0$ such that
\begin{align*}
    q_{\best{\pol}}(s_N, \best{\pol}) > q_{\best{\pol}}(s_N, a) + \epsilon
\end{align*}
for all actions $a \neq \best{\pol}(s_N)$.
Therefore, the definition of convergence indicates that there exists a time step $K_N$ such that for all $k \geq K_N$
\begin{align*}
    q_k(s_N, \best{\pol}) 
    \geq q_{\best{\pol}}(s_N, \best{\pol}) - \epsilon/2 
    > q_{\best{\pol}}(s_N, a) + \epsilon/2 
    \geq q_k(s_N, a) 
\end{align*}
for all actions $a \neq \best{\pol}(s_N)$.
Consequently, for all $k \geq K_N$ it holds that
\begin{align*}
    \pol_k(s_N) = \best{\pol}(s_N) .
\end{align*}

Further, consider any state $s_n$ and suppose there exists a time step $K_{n+1}$ such that for all steps $k \geq K_{n+1}$
\begin{align*}
    \pol_k(s_n') = \best{\pol}(s_n')
\end{align*}
for all $n'>n$.
The expected return of $(s_n, \best{\pol})$ for every episode after $K_{n+1}$ is then $q_{\best{\pol}}(s_n, \best{\pol})$. As a result, the estimate is updated infinitely often as
\begin{align*}
    q_\kpo(s_n, \best{\pol})
    =
    q_k(s_n, \best{\pol})
    + \lrate_k \update_k ( q_{\best{\pol}}(s_n, \best{\pol}) - q_k(s_n, \best{\pol}) + w_k(s_n, \best{\pol}) ) 
\end{align*}
and, therefore, it converges to $q_{\best{\pol}}(s_n, \best{\pol})$ with probability one.
As explained above, all other actions are bounded by $q_{\best{\pol}}(s_n, a)$.
So there exists a time step $K_n \geq K_{n+1}$ such that for all $k \geq K_n$
\begin{align*}
    q_k(s_n, \best{\pol}) 
    > q_k(s_n, a) 
\end{align*}
for all actions $a \neq \best{\pol}(s_n)$.
Consequently, for all $k \geq K_n$ it holds that
\begin{align*}
    \pol_k(s_{n'}) = \best{\pol}(s_{n'})
\end{align*}
for all $n' \geq n$.

This induction demonstrates that the policy $\pol_k$ matches the policy $\best{\pol}$ for all $k \geq K_1$. For every episode, the expected return of every state-action pair $(s,a)$ is then $q_{\best{\pol}}(s,a)$. As a result, for all $k \geq K_1$ the estimate is updated as
\begin{align*}
    q_\kpo
    =
    q_k
    + \lrate_k \update_k ( q_{\best{\pol}} - q_k + w_k ) 
\end{align*}
and, therefore, it converges to $q_{\best{\pol}}$ with probability one.
\end{proof}

% The proof of \autoref{thm: optimal feedforward} hinges on 
% % the fact that if the optimal policy is feedforward, there exists a state $s$ such that for all policies 
% % \begin{align}
% %     q_\pol(s,\best{\pol}) &= q_\best{\pol}(s,\best{\pol})
% %     \\
% %     q_\pol(s,\best{\pol}) &> q_\pol(s,a)
% % \end{align}
% the acyclic graph induced by the feedforward optimal policy.
% However, this property is difficult to test in practice, especially if the MDP is unknown.

% \begin{remark}
%     can apply same argument on groups of states
    
%     suppose the optimal policy splits the MDP into 
%     feedforward branches, i.e. set of states with zero revisit probability
%     and looping branches set of states with non-zero revisit probability

%     argument shows that if the end of a feedforward branch has converged, then the entire feedforward branch converges

%     thus, the open problem is really about the convergence of MC-O-PI in the states where the optimal policy loops
%     only need to study MDP where for each state there is a non-zero probability of looping back to it
% \end{remark}

\vspace{-0.2cm}


%%%%%%%%%%%%%%%%%%%%%%%%%%%%%%%%%%%%%%%%%
%%%%%%%%%%%%%%%%%%%%%%%%%%%%%%%%%%%%%%%%%
%%%%%%%%%%%%%%%%%%%%%%%%%%%%%%%%%%%%%%%%%
\section{Conclusion}

% The stochastic difference inclusion \eqref{eq: complete difference inclusion} underlying 
MC-O-PI has favourable properties as solutions converge to a bounded set and perturbations caused by episode simulations and asynchronous updates cannot prevent convergence.
However, its asymptotic behaviour remains complex. Under the modified Robbins-Monro conditions alone, solutions may converge to a suboptimal fixed point. With extra assumptions on the update distributions or on the optimal policy, they almost surely converge to $\qopt$.
Building on these insights, this thesis develops more general convergence proofs and more realistic counterexamples.


% need to find convergence or counterexamples when
% - update distributions are not uniform over all state-action pairs
% - optimal policy is not feedforward

% explain exactly what we do
% based on what has already been done


\clearpage